\chapter{The module \texorpdfstring{\texttt{(aux kanren lambda)}}{(aux kanren lambda)}}\label{ch:module}

% ------------------------------------------------------------------------------------------------
\section*{What and why}

\Cref{ch:micro} took \micro{} apart and ended, in \cref{sec:micro-ext}, with the list of points an extension
cannot reach from outside: the state has no free slot, and the unifier, the walkers and the reifier are closed.
\Cref{ch:design} opened exactly those points with a small, λ-agnostic protocol of hooks and explained why the
proposal is a client module of that protocol rather than a native engine. This chapter is about the client
itself: the module \scm{(aux kanren lambda)}, which brings into \micro{} the three ingredients of Part~I,
\begin{itemize}
  \item λ-tree syntax (\cref{ch:lambdatree}): object-level binders are meta-level abstractions;
  \item higher-order pattern unification, $\Llam$ (\cref{ch:hounif}), with universe levels in place of the
        mixed prefix and with non-patterns suspended (\enquote{dynamic $\Llam$});
  \item universal and hypothetical goals, \rulename{generic} and \rulename{augment}, over hereditary Harrop formulas (\cref{ch:hh}).
\end{itemize}

The chapter is the formal heart of the proposal. It states what the terms are (\cref{sec:module-terms}),
what the module keeps in the state (\cref{sec:module-state}), the unification rules with their side
conditions and the termination argument (\cref{sec:module-unif}), the rules for \scm{∀°} and \scm{⊃°} and
why levels make them sound (\cref{sec:module-hh}), how micro's disequality, absence and type constraints
read λ-terms (\cref{sec:module-stores}), how answers are printed (\cref{sec:module-reify}), and, as a
reference, the exported API (\cref{sec:module-api}). The curated book examples are the subject of
\cref{ch:examples}; the known limits, the costs and the merge plan are in \cref{ch:roadmap}.

{\emergencystretch=2em
The module is 1255 lines, 885 of them neither blank nor comment. In the repository it lives, with its
tests and examples, under \texttt{doc/\allowbreak hol-kanren/\allowbreak prototype/}, and after merging as
\texttt{src/\allowbreak aux.kanren.lambda.scm}, with the suites \texttt{src/\allowbreak test/\allowbreak microkanren-lambda.scm} (83 cases, 320
assertions) and \texttt{src/\allowbreak test/\allowbreak microkanren-lambda-limits.scm} (6 cases). Every Scheme snippet below was run
against the module and micro with the protocol applied; its output is pasted verbatim under it. Before
the imports below work, the edited micro must be built: \cref{sec:roadmap-run} shows how, without
installing anything, and \cref{sec:roadmap-merge} how to merge it. The snippets
import\par}
\begin{scheme}
(import scheme (chicken base) (chicken pretty-print)
        (chicken condition) (chicken memory representation)
        (aux base) (aux stream) (aux kanren micro) (aux kanren micro show)
        (aux kanren lambda))
\end{scheme}
Several of them use a \emph{probe}: a goal that prints part of the state and succeeds. A goal of \micro{} is a
procedure from a state to a stream of states (\cref{def:micro-goal}), so a probe is three lines; each snippet
shows the one it uses.

% ------------------------------------------------------------------------------------------------
\section{Terms}\label{sec:module-terms}

\subsection{Four records}

First-order structure stays Scheme data, exactly as in \micro: pairs, vectors, foreign records and atoms
(\cref{def:micro-term}). The module adds four record types, defined at the top of the file:
\begin{scheme}
(define-record μkanren-lam hint body)        ; λ.body, locally nameless
(define-record μkanren-app head args)        ; head applied to a list of arguments
(define-record μkanren-eigen name level id)  ; a scoped constant of GENERIC
(define-record μkanren-bvar index)           ; a de Bruijn index, only inside bodies
\end{scheme}

\begin{definition}[λ-terms of the module]\label{def:module-terms}
Extend the terms of \micro{} (\cref{def:micro-term}) with
\[
  t \;::=\; \dots \;\mid\; c^{\ell} \;\mid\; \#i \;\mid\; \lambda^{h}.\,t \;\mid\; t\;@\;\langle t_1,\dots,t_n\rangle
\]
where $c^\ell$ is an eigenvariable of level $\ell\in\mathbb N\cup\{\infty\}$ (\scm{μkanren-eigen}), $\#i$ a
bound variable with de Bruijn index $i\ge 0$ (\scm{μkanren-bvar}), $\lambda^h.t$ an abstraction whose binder
carries the printing hint $h$, a symbol (\scm{μkanren-lam}), and $t\,@\,\bar t$ an application with $n\ge 1$
(\scm{μkanren-app}). A term is \emph{locally closed} when every $\#i$ in it lies below at least $i+1$
abstractions. Only locally closed terms are ever bound to variables or stored in constraints.
\end{definition}

The representation is \emph{locally nameless} \cite{Chargueraud2012}: bound variables are de Bruijn indices
\cite{deBruijn1972}, free ones are names (logic variables and eigenvariables). Three design choices follow
from it.
\begin{itemize}
  \item \textbf{The hint is only for printing}, and it is an \emph{interned} symbol. The binder's identity is
    its position. Two terms built alike are therefore \scm{equal?}, so λ-terms can be grouped by
    \scm{groupby°} and compared by the host.
  \item \textbf{An eigenvariable's identity is its \texttt{id}}, an uninterned symbol, compared with \scm{eq?}.
    Its \texttt{name} is a hint and its \texttt{level} is the universe it lives in (\cref{sec:module-state}).
    Level $\infty$ (\scm{+inf.0}) marks the names that the unifier opens when it descends under binders and
    the placeholders that \scm{ƛ} uses while building a body.
  \item \textbf{The first 64 bound variables are preallocated} in a vector, so the common indices allocate
    nothing.
\end{itemize}
The de Bruijn operations (shift, instantiate, close) never walk the substitution: a term bound to a variable
is locally closed, so a binding never contributes a loose index. Instantiation shifts the substituted term
only when it has loose indices itself, which is decided once per instantiation.

\subsection{The constructors \texorpdfstring{\texttt{ƛ}}{ƛ} and \texorpdfstring{\texttt{@}}{@}}

The user never builds these records by hand. Two constructors give the surface syntax of λ-tree terms:
\begin{scheme}
(define-syntax ƛ
  (syntax-rules ()
    ((_ () body) body)
    ((_ (x y ...) body) (μkanren-abstract 'x (λ (x) (ƛ (y ...) body))))))

; λ.(f p), with the placeholder p closed into #0
(define (μkanren-abstract name f)
  (let1 (p (μkanren-eigen* name +inf.0))
    (μkanren-lam/η name (μkanren-close (f p) p))))

(define (@ f . args) (μkanren-beta f args))   ; β at once on an abstraction
\end{scheme}
\scm{(ƛ (x) body)} binds the \emph{Scheme} variable \scm{x} to a fresh placeholder eigenvariable while
\scm{body} is evaluated, then closes that placeholder into the index \#0. The object binder is thus a Scheme
binder: shadowing, lexical scope and quasiquotation all behave as a Scheme programmer expects, which is the
point of λ-tree syntax \book{177}. Two invariants are enforced at construction.
\begin{itemize}
  \item \textbf{η-short.} \scm{μkanren-lam/η} builds $\lambda x.\,(h\,\bar a\,x)$ as $h\,\bar a$ when $x$ is not
    free in $h\,\bar a$. The unifier builds its abstractions the same way. This matters because micro's own
    occurs check runs before any hook can see an equation: $\lambda x.\,(F\,x) = F$ holds by η, but in
    η-long form micro would see $F$ inside the value it binds $F$ to and fail; built η-short, the left side
    \emph{is} $F$ and the equation is trivial.
  \item \textbf{β at construction.} \scm{@} applied to an abstraction β-reduces at once, and nested
    applications are flattened: \scm{(@ (@ f a) b)} is \scm{(@ f a b)}.
\end{itemize}

\begin{example}[Construction and printing]\label{ex:module-construct}
Every record has a printer that shows its term as data. The raw record of $\lambda x.\,(f\ x)$ shows the hint
\texttt{x} and the index in the body:
\begin{scheme}
(ƛ (x y) (cons x y))
(record->vector (ƛ (x) `(f ,x)))     ; from (chicken memory representation)
(equal? (ƛ (x) `(f ,x)) (ƛ (x) `(f ,x)))
(equal? (ƛ (x) x) (ƛ (y) y))
(μkanren-α=? (ƛ (x) x) (ƛ (y) y))
(@ (ƛ (x) `(g ,x ,x)) 'a)
(ƛ (x) (@ 'g x))
(ƛ (x) `(g ,x))
(@ 'g 'a)
(equal? (@ 'g 'a) '(g a))
(ƛ (a) (ƛ (a) a))
\end{scheme}
\begin{answers}
  => (ƛ (x y) (x . y))
  => #(aux.kanren.lambda#μkanren-lam x (f #<aux.kanren.lambda#μkanren-bvar>))
  => #t
  => #f
  => #t
  => (g a a)
  => g
  => (ƛ (x) (g x))
  => (@ g a)
  => #f
  => (ƛ (a a1) a1)
\end{answers}
Line by line: nested binders print merged, \scm{(ƛ (x y) ...)}; the body of a \texttt{μkanren-lam} holds an
index, not the name; terms built alike are \scm{equal?}; α-variants with \emph{different} hints are not
\scm{equal?} but are α-equal; \scm{@} reduced the redex; $\lambda x.\,(g\,x)$ written with \scm{@} is the
η-redex of $g$ and was built as \texttt{g}; written with a quasiquoted list it is no η-redex (the list
\scm{(g x)} is data, not an application), so it stays an abstraction; the λ-tree application of a constant
prints with \texttt{@} and is not the list \scm{(g a)}; shadowing is lexical, and the printer renames the
inner binder apart (\cref{sec:module-reify}).
\end{example}

\subsection{Three kinds of application}

An application is classified by its head \emph{after walking it} in the current substitution:
\begin{definition}[Flexible, β-redex, rigid; head normal form]\label{def:module-hnf}
An application $h\,@\,\bar t$ is \emph{flexible} if $h$ walks to an unbound logic variable, a
\emph{β-redex} if $h$ walks to an abstraction, and \emph{rigid} otherwise (a constant, an eigenvariable, a
bound variable). The \emph{head normal form} $\mathrm{hnf}(t)$ (\scm{μkanren-hnorm}) walks $t$, flattens
nested applications, contracts head β-redexes, and contracts an abstraction whose body has \emph{become} an
η-redex through a binding (e.g.\ $\lambda x.\,F\,(G\,x)$ once $G:=\lambda z.z$).
\end{definition}
Under the $\Llam$ restriction every contraction hnf performs is a $\beta_0$ renaming
(\cref{def:lambdatree-beta0}). On instantiated heads \scm{@} gives full β, which is what the book calls
dynamic $\Llam$ \book{225--226}.
\begin{scheme}
(μkanren-hnorm (@ (ƛ (x) (@ x 'b)) (ƛ (y) `(g ,y ,y))) μkanren-state-empty)
(°->list #f (fresh° (q F) (=° q (@ F 'a))))
(°->list/datum (fresh° (q F) (=° q (list F (@ F 'a))) (=° F (ƛ (z) `(h ,z)))))
\end{scheme}
\begin{answers}
(g b b)
((λ (α) (@ α 'a)))
(((ƛ (z) (h z)) (h a)))
\end{answers}
The second answer is a flexible term, printed as \scm{(@ α 'a)}; once $F$ is bound (third line) the same
term is a β-redex and its value is \scm{(h a)}.

\subsection{Scheme lists versus λ-tree application}

The object language of the book's programs is written with Scheme lists, as \micro{} always did:
$\mathtt{abs}\ M$ is \scm{`(abs ,M)}, $\mathtt{app}\ M\ N$ is \scm{`(app ,M ,N)}, with $M$ a meta-level
abstraction for a binder. The λ-tree application \scm{(@ 'g 'a)} and the list \scm{'(g a)} are distinct terms
and never unify:
\begin{scheme}
(°->list #f (fresh° (q) (=° (@ 'g 'a) '(g a))))
\end{scheme}
\begin{answers}
  => ()
\end{answers}
This is deliberate. One of the three prototypes identified the two (a spine notation for lists) and lost
commutativity of conjunction: with η, the same pair of goals gave different answers in the two orders. Keeping
lists as data and \scm{@} as the only application keeps \micro's first-order unification untouched for lists,
and confines the higher-order rules to terms that contain a λ-record.

\begin{pitfall}[What is and is not an η-redex]
η-contraction applies to λ-tree applications only. \scm{(ƛ (x) (@ 'g x))} \emph{is} \texttt{g}, but
\scm{(ƛ (x) `(g ,x))} is an abstraction over a two-element list, and stays one. The module is untyped, so it
cannot η-\emph{expand} a term against its type: η-expansion happens only on demand, when one side of an
equation is an abstraction (rule \rulename{η} below), as in Miller's untyped treatment
\cite[§9.3]{Miller1991}. A rigid list such as \scm{'(g a)} therefore equals an abstraction only when the
abstraction's body equals the λ-tree application \scm{(@ '(g a) x)} of the list to the bound variable (and
\scm{ƛ} already builds such an η-redex as the list itself); never when the body is a list such as
\scm{`(g a ,x)}.
\end{pitfall}

% ------------------------------------------------------------------------------------------------
\section{The extension's state}\label{sec:module-state}

\subsection{The record in \texorpdfstring{\texttt{X}}{X}}

The protocol gives \micro's state a seventh slot \texttt{X} that micro copies but never reads
(\cref{ch:design}). The module keeps one record in it:
\begin{scheme}
(define-record μkanren-hh depth top L Ln Γ E N)
(define μkanren-hh0 (make-μkanren-hh 0 0 empty/sbral 0 '() '() '()))
(define (μkanren-hh s) (let1 (x (μkanren-state-X s)) (if (null? x) μkanren-hh0 x)))
\end{scheme}
\texttt{X} = \scm{'()}, the value in \scm{μkanren-state-empty}, reads as \enquote{everything empty}, so the
initial state of every runner works unchanged.

\begin{definition}[The λ state]\label{def:module-state}
The record $\langle d,\ \mathit{top},\ L,\ L_n,\ \Gamma,\ E,\ N\rangle$ holds:
\begin{description}[style=nextline,leftmargin=2.2em]
  \item[$d$ (\texttt{depth})] the current universe level, raised by \scm{∀°} for the extent of its body;
  \item[$\mathit{top}$] the highest level reached so far in this branch;
  \item[$L,\,L_n$] the levels of the logic variables: $L$ is a skew-binary random-access list (\aux's
    \texttt{sbral}) indexed like the substitution $S$, holding the levels of the variables with index
    $<L_n$. The level of variable $i$ is
    \[
      \ell(i) \;=\; \begin{cases} 0 & i<0 \text{ (a reified variable)}\\
                                   L[L_n-1-i] & 0\le i<L_n\\
                                   d & i \ge L_n \end{cases}
    \]
    so a variable created since $L$ was last extended has the current depth. $L$ is padded with $d$ exactly
    when the depth changes and whenever the unifier creates a variable of another level; lookup is
    $O(\log n)$;
  \item[$\Gamma$] the assumed clauses: an alist from a relation to its clauses, most recent first, one entry
    per distinct relation;
  \item[$E$] the suspended constraints: equations \scm{(= u v)} outside $\Llam$ (or bindings whose scope check
    a non-pattern decides), types \scm{(tag tag u)} on flexible terms, and absentos \scm{(absent d u)} whose
    datum outlived its eigenvariable's scope;
  \item[$N$] the disequalities that involve λ-terms or eigenvariables, as \scm{(reading u . v)}. micro's D never
    holds one (\cref{sec:module-stores}).
\end{description}
\end{definition}

Levels are the module's whole account of quantifier order. A logic variable created at depth $d$ stands for
an $\exists$ placed after the first $d$ universals of its branch; an eigenvariable created by the $(d{+}1)$-th
\scm{∀°} has level $d+1$. The binding conditions of Nadathur, Jayaraman and Kwon \cite{NadathurJayaramanKwon1995}
(\cref{def:rules-binding,def:lkanren-ui}) are then stated on levels: a variable of level $\ell$ may be bound
only to a term whose eigenvariables have level $\le\ell$, and the variables of level $>\ell$ in that term are
\emph{lowered} to $\ell$.

\subsection{Copied, never mutated}

The record is immutable. Every change builds a new one with \scm{μkanren-hh-with}, a keyword-argument copy, and
installs it with \scm{(μkanren-state-with s X: x)}. Since a state is a value and \scm{or°} gives each branch the
state it was given, two branches can never see each other's depth, Γ or suspended constraints. What needs
care is the \emph{end} of a scope: \scm{∀°} and \scm{⊃°} change the record for their body only, and each
answer of the body must give the change back. Both do so on every answer state (\cref{sec:module-hh}). The
following probe shows depth and levels through two nested \scm{∀°}:
\begin{scheme}
(define ((probe label . vars) s)
  (display label) (display ": depth ") (display (μkanren-λ-depth s))
  (display ", levels ") (write (map (λ (v) (μkanren-λ-level v s)) vars))
  (display ", E ") (write (μkanren-λ-delayed s))
  (display ", N ") (write (μkanren-λ-disequalities s)) (newline)
  (✓° s))
(°->list #f
  (fresh° (q)
    (probe "start" q)
    (∀° (x)
      (fresh° (Y)
        (probe "inside ∀x" q Y)
        (∀° (z) (fresh° (W) (probe "inside ∀x∀z" q Y W)))
        (probe "back in ∀x" q Y)
        (=° (@ q x) (@ Y x x))
        (probe "after q x = Y x x" q Y)))
    (probe "after ∀x" q)))
(°->list #f (fresh° (q) (or° (∀° (x) (probe "branch 1" q) (=° q 1))
                              (and° (probe "branch 2" q) (=° q 2)))))
\end{scheme}
\begin{answers}
start: depth 0, levels (0), E (), N ()
inside ∀x: depth 1, levels (0 1), E (), N ()
inside ∀x∀z: depth 2, levels (0 1 2), E (), N ()
back in ∀x: depth 1, levels (0 1), E (), N ()
after q x = Y x x: depth 1, levels (0 1), E (), N ()
after ∀x: depth 0, levels (0), E (), N ()
((λ (α) (ƛ (x) (@ α x x x))))
branch 1: depth 1, levels (0), E (), N ()
branch 2: depth 0, levels (0), E (), N ()
((λ () 1) (λ () 2))
\end{answers}

\begin{tutorial}[Reading the probe]
$q$ is created at depth~0, $Y$ inside \scm{∀° (x)} at depth~1, $W$ at depth~2; after each \scm{∀°} the depth
is back where it was. Written as a mixed prefix, the goal is $\exists q\,\forall x\,\exists Y\,\forall z\,\exists W$,
and each level is the number of universals to the left of the quantifier: $q{:}0$, $x{:}1$, $Y{:}1$, $z{:}2$,
$W{:}2$. The equation $q\,x \doteq Y\,x\,x$ is solved inside the scope, and its answer shows the
levels at work. $q$ has level~0 and its argument $x$ (level~1) is a pattern for it, so the unifier eliminates
$q$: $q := \lambda x.\,t'$ where $t'$ is $Y\,x\,x$ with $x$ turned into the bound variable. But $Y$ has
level~1 and $q$ level~0: $q$'s value may not mention a variable that could see $x$ unless $x$ is passed
explicitly. $Y\,x\,x$ is no pattern ($x$ is repeated), but both its arguments are $q$'s own argument $x$,
so none of them has to be dropped. $Y$ is therefore \emph{lowered} and \emph{raised} at once, $Y := \lambda a\,b.\,Y'\,a\,b\,x$ with
$\ell(Y')=0$, and $q := \lambda x.\,Y'\,x\,x\,x$. The answer \scm{(ƛ (x) (@ α x x x))} is that binding; the
rules responsible are \rulename{FR} and case~(6c) of variable elimination, \cref{def:module-elim}. The last
two lines show the branches: \scm{∀°} raised the depth in branch~1 only.
\end{tutorial}

The module's hooks keep the two protocol invariants of \cref{def:design-invariants}; in particular, every
function of the module that may leave the state unchanged is written to return it \scm{eq?}.

% ------------------------------------------------------------------------------------------------
\section{Unification}\label{sec:module-unif}

\subsection{Where the rules run}

\micro's unifier \scm{μkanren-state-unify} still does all first-order work: variables, pairs, vectors, foreign
records and atoms. Through the protocol's \texttt{unify} hook it hands any pair with a walked λ-record on one
side, neither side a variable, to \scm{μkanren-λ-unify}, which applies the rules of \cref{fig:module-rules}.
The rules call back into \scm{μkanren-state-unify} for subterms, so the two unifiers interleave freely. After
\scm{=°} has made its bindings, the \texttt{post} hook \scm{μkanren-λ-post} runs before micro verifies D, A
and T:
\begin{enumerate}
  \item \emph{settle}: the scope check \textbf{Check} of every new binding, then the wake-up of $E$, in rounds
    while a round binds (\cref{sec:module-wake});
  \item \emph{migrate}: the prefixes of D that now hold a λ-term or an eigenvariable move to N, and the types of
    A on variables just bound to λ-terms move to the \texttt{tag} hook (\cref{sec:module-stores});
  \item \emph{verify N}: every λ-disequality is tried again.
\end{enumerate}

\subsection{Notation}

Write $\ell(X)$ for the level of a logic variable and $\ell(c)$ for that of an eigenvariable; the names the
unifier opens under binders have $\ell(c)=\infty$. A \emph{name} is an eigenvariable or, inside variable
elimination, a bound variable of the term being walked.
\begin{definition}[Pattern, read off levels]\label{def:module-pattern}
A list of arguments $\bar a = a_1\dots a_n$ is a \emph{pattern for} a variable $F$ if every $\mathrm{hnf}(a_i)$
is an eigenvariable, the $a_i$ are pairwise distinct, and $\ell(a_i) > \ell(F)$ for every $i$
(\scm{μkanren-pattern?}). A flexible term $F\,\bar a$ is a pattern when $\bar a$ is a pattern for $F$.
\end{definition}
For example, inside \scm{(∀° (x y) ...)} with $\ell(F)=0$, $\ell(x)=1$, $\ell(y)=2$, the lists $\langle x,y\rangle$
and $\langle y\rangle$ are patterns for $F$, while $\langle x,x\rangle$ (repeated), $\langle \mathtt{a}\rangle$ (a
constant) and $\langle Y\rangle$ (a variable) are not. The same definition, with a bound variable of the
term being walked counted as an eigenvariable of level $\infty$, is used inside variable elimination (case (6b)
below).
This is Miller's $\Llam$ condition (\cref{def:hounif-llambda}): the arguments are distinct variables
\enquote{quantified inside the scope of $F$}. With a prefix, \enquote{inside the scope} means \enquote{to the
right of $F$}; with levels it means \enquote{of a level above $F$'s}, which is how Nadathur and Linnell read the
condition without a prefix \cite[§4]{NadathurLinnell2005}. An eigenvariable $F$ could already see is not a
pattern argument: in $\forall x\,\exists F.\ F\,x = \dots$ the term $F\,x$ is no pattern, since $F$ may use
$x$ directly and the equation has several incomparable solutions.

Throughout, $\lambda.b$ is an abstraction with body $b$, $b[c/0]$ its body with \#0 instantiated to $c$,
$h\,\bar u$ an application, and $F,G,H,X,Y$ logic variables. Both sides are put in hnf before a rule is
chosen.

\subsection{The rules}

The rules of \cref{fig:module-rules} are transformations, read top-down as in \cref{fig:rules-fo,fig:rules-pattern}: the equation above the
line is replaced by the equations, bindings or suspension below it. They are the rules of \cref{fig:rules-pattern}
with the prefix erased into levels; \cref{sec:hounif-algo} gives the same rules under other names, recalled
in parentheses. \rulename{ξ} is \rulename{abs} (\rulename{Abs}), \rulename{η} is \rulename{eta}
(\rulename{Eta}), \rulename{Rigid} is \rulename{rr-dec} and \rulename{rr-clash} (\rulename{Dec} and
\rulename{Clash}), \rulename{FF-same} is \rulename{ff-same} (the same-head case of \rulename{FlexFlex}),
\rulename{Delay} is \rulename{delay}. \rulename{FR} covers \rulename{elim} (\rulename{Elim}) with its pruning,
\rulename{occurs} and \rulename{scope}, and also \rulename{ff-diff} (\rulename{FlexFlex} with distinct heads): instead of binding both heads to a new $H$,
it eliminates $F$ and prunes $G$ only where needed, which gives the same mgu up to the name of $H$
(\cref{ex:module-ff}). \rulename{Var} and \rulename{Occurs} are \rulename{elim} and \rulename{occurs} for a
variable with no arguments, a case that micro's first-order unifier meets first, and they are where the
eager scope check, \textbf{Check}, enters; \rulename{FF-one} orients a flexible–flexible pair as
\rulename{swap} orients a flexible–rigid one.

\begin{figure*}[!t]
\begin{mathpar}
\inferrule{X \doteq t}{X := t\ \ \text{(micro, occurs check)}\quad\textbf{Check}(X := t)}\;\rulename{Var}
\and
\inferrule{\lambda.b \doteq \lambda.b'}{b[c/0] \doteq b'[c/0]\quad \textbf{Check}\ \text{its bindings}}\;\rulename{ξ}
\and
\inferrule{\lambda.b \doteq t}{b[c/0] \doteq t\;@\;c\quad \textbf{Check}\ \text{its bindings}}\;\rulename{η}
\\
\inferrule{h\,\bar u \doteq h'\,\bar v}{h \doteq h' \quad \bar u \doteq \bar v}\;\rulename{Rigid}
\and
\inferrule{F\,\bar a \doteq F\,\bar b}{F := \lambda\bar x.\,H\,\langle x_i \mid a_i = b_i\rangle}\;\rulename{FF-same}
\\
\inferrule{F\,\bar a \doteq t}{F := \lambda\bar x.\,t'\qquad t' = \mathrm{Elim}(\ell(F),F,\bar a,t)}\;\rulename{FR}
\and
\inferrule{F\,\bar a \doteq G\,\bar b}{G\,\bar b \doteq F\,\bar a\ \ \text{by \rulename{FR}}}\;\rulename{FF-one}
\\
\inferrule{u \doteq v}{E := E \cup \{\lambda L.\,u \doteq \lambda L.\,v\}}\;\rulename{Delay}
\and
\inferrule{X \doteq t \\ X \text{ occurs in } t}{\top\ \text{if } \mathrm{hnf}(t)=X;\ \text{else } X := \mathrm{Elim}(\ell(X),X,\langle\rangle,t)}\;\rulename{Occurs}
\end{mathpar}
\caption{The unification rules of \texttt{μkanren-λ-unify}, applied after hnf of both sides; above the line
the equation, below what is done with it. In \rulename{ξ} and \rulename{η}, $c$ is a new eigenvariable of level
$\infty$. In \rulename{FF-same} and \rulename{FR}, $\bar x$ are new binders, one per argument $a_i$, and Elim
replaces each $a_i$ in $t$ by the bound variable of $x_i$. Variable elimination, Elim, is
\cref{def:module-elim}; \textbf{Check} is \cref{def:module-check}. Each rule may also answer $\bot$ or suspend,
as its side conditions say.}\label{fig:module-rules}
\end{figure*}

The side conditions, in the order the code tests them (\scm{μkanren-λ-unify} and the functions it calls):
\begin{description}[style=nextline,leftmargin=5.6em]
  \item[\rulename{Var}] One side is an unbound variable. micro binds it, with its occurs check; the scope check
    \textbf{Check} runs afterwards, in the \texttt{post} hook. Variable–variable pairs are micro's as well.
  \item[\rulename{ξ}] Both sides are abstractions. $c$ is a new eigenvariable of level $\infty$, so no variable
    can capture it except through its own pattern arguments. Because $c$ is local to this very unification,
    the bindings made under it are checked at once (\scm{μkanren-scoped}), not at the end of \scm{=°}.
  \item[\rulename{η}] One side is an abstraction and the other is not. The other side is applied to the new
    $c$. This is the only η-expansion the module performs.
  \item[\rulename{Rigid}] Both sides are rigid applications with the same number of arguments: heads and
    argument lists are unified by micro. Two eigenvariables unify iff they are the same (\scm{eq?} ids), two
    bound variables iff their indices agree. Any other pair with a λ-term on one side fails: different rigid
    heads, different arities (untyped), a λ-record against a pair or an atom.
  \item[\rulename{FF-same}] Same head on both sides, both argument lists patterns. Equal lists: success with
    no binding. Different lengths: $\bot$. Otherwise $F$ keeps the positions where the arguments agree, with a
    new $H$ of level $\ell(F)$ \book{222}. If $\bar b$ is no pattern, \rulename{Delay}.
  \item[\rulename{FR}] One side is a pattern $F\,\bar a$; the other side $t$ is rigid, or flexible with another
    head (\enquote{FF-diff}: variable elimination prunes $G$ as needed \book{222}). If Elim fails the rule
    fails; if Elim suspends, \rulename{Delay}. The binding $\lambda\bar x.\,t'$ is built η-short, its binders
    named after the eigenvariables $\bar a$.
  \item[\rulename{FF-one}] Two flexible terms, only the right one a pattern: \rulename{FR} on $G$.
  \item[\rulename{Delay}] No side is a usable pattern. The equation is pushed on $E$, closed over the names $L$
    that the unifier had opened so far ($\lambda L.u \doteq \lambda L.v$, so no $\infty$-level name escapes the
    unification that created it), \emph{unless} $u$ and $v$ are already α-equal, in which case it is dropped.
  \item[\rulename{Occurs}] Not a rule of \scm{μkanren-λ-unify} but of the protocol's \texttt{occurs} hook: micro's
    occurs check failed on $X := t$ and $t$ holds a λ-term. If $\mathrm{hnf}(t)$ is $X$ itself the equation
    holds ($q = F\,q$ with $F := \lambda z.z$); otherwise $X$ is eliminated with no arguments, which may fail,
    succeed, or suspend ($F = \lambda x.\,G\,(F\,x)$ holds with $G := \lambda z.z$). A first-order cycle,
    $X = f\,X$, holds no λ-term and fails in micro as before.
\end{description}

\subsection{Variable elimination}

\rulename{FR} is the book's variable elimination \book{223} (\cref{def:hounif-elim}) with levels in place of the
prefix. It walks the right-hand side once and builds the body of $F$'s binding, pruning, lowering and raising the
variables it meets on the way. Compared with \cref{def:hounif-elim}, three things are added: \emph{lowering}
with delayed raising (case~(4)), where \cref{def:hounif-elim} assumes a prefix already raised into
$\forall\exists$ form (the eager raising of \cref{ex:hounif-8k}); the treatment of \emph{non-patterns} inside the right-hand side (cases (6c) and (7)), which the book never meets since
it stays in $\Llam$; and a finer \emph{occurs} test (case (6a)), since an occurrence of $F$ is fatal only when no
instance can make it smaller. The contexts below exist to decide these three cases.

\begin{definition}[Variable elimination]\label{def:module-elim}
$\mathrm{Elim}(\ell, F, \bar a, t)$, with $\bar a$ a pattern for $F$ and $\ell=\ell(F)$ (or $F$ absent and
$\bar a$ empty, for \textbf{Check}), walks $\mathrm{hnf}(t)$ at a local depth $k$ (the binders crossed inside
$t$) and in a \emph{context}: \emph{strict} (\texttt{top}: below binders only; \texttt{rigid}: below a first rigid
constructor that is an application headed by a bound variable or by an $a_i$; \texttt{firm}: below a first rigid
constructor that is a constant, an eigenvariable not in $\bar a$, a pair, a vector or a record, a root no instance
of $F$ can project onto) or \emph{soft} (inside an argument of a flexible term that is no pattern, or
\texttt{weak}: below a flexible head inside such an argument). It answers a term $t'$ with new bindings,
$\bot$, or \emph{suspend}. By cases on the walked subterm:
\begin{enumerate}[label=(\arabic*)]
  \item \emph{Eigenvariable $c$}: if $c = a_i$, the bound variable of $a_i$; else if $\ell(c)\le\ell$, $c$
    itself; else $\bot$ (scope failure, \cref{ex:hounif-8j}). In a weak context the last case is trouble, not
    $\bot$.
  \item \emph{Bound variable}: itself.
  \item \emph{Abstraction} $\lambda.b$: $\lambda.\mathrm{Elim}(b)$ at depth $k+1$.
  \item \emph{Unbound variable $Y$}: if $Y=F$, $\bot$ when $\bar a$ is empty (the first-order occurs check) and
    suspend otherwise; if $\ell(Y)\le\ell$, $Y$; in a soft context, trouble; otherwise \emph{lower and raise}:
    $Y := Y'\,\bar a_Y$ with $\ell(Y')=\ell$ and $\bar a_Y = \langle a_i \mid \ell(a_i)\le\ell(Y)\rangle$, the
    arguments of $F$ that $Y$ could see; the result is $Y'$ applied to their bound variables
    (\enquote{delayed raising}, \cref{ex:hounif-8k}).
  \item \emph{Rigid application, pair, vector, record}: homomorphic. A \texttt{top} context becomes
    \texttt{rigid} below an application whose head is a bound variable or in $\bar a$, and \texttt{firm} below
    any other head, a pair, a vector or a record; any other context is passed on unchanged.
  \item \emph{Flexible $G\,\bar b$ in a strict context}:
    \begin{enumerate}[label=(\alph*)]
      \item if $G=F$: $\bot$ (occurs failure) when the context is not \texttt{top}, no $b_j$ is
        \emph{functional} (an abstraction, a variable or a flexible term, which may become an abstraction), and
        either the context is \texttt{firm} or $|\bar b| = |\bar a| + k_{\mathrm{top}}$, with $k_{\mathrm{top}}$
        the binders on top of $t$; otherwise suspend;
      \item if $\bar b$ is a pattern for $G$ (each $b_j$ a local bound variable or an eigenvariable of level
        $>\ell(G)$, pairwise distinct): keep $b_j$ iff it is local, or in $\bar a$, or $\ell(b_j)\le\ell$; let
        $\mathit{extra}=\langle a_i \mid \ell(a_i)\le\ell(G)\rangle$ if $\ell(G)>\ell$, else empty. If every
        $b_j$ is kept, $\mathit{extra}$ is empty and $\ell(G)\le\ell$, $G\,\bar b$ is unchanged. Otherwise
        \emph{prune}: $G := \lambda\bar y.\,G'\,\bar y|_{\mathit{keep}}\,\mathit{extra}$ with
        $\ell(G')=\min(\ell,\ell(G))$, and the result is $G'$ applied to the images of the kept arguments and of
        $\mathit{extra}$ (\cref{ex:hounif-8h});
      \item if $\bar b$ is no pattern: each $b_j$ is walked on its own in a soft context, binding nothing; it is
        \emph{kept}, as the walk rebuilt it (its $a_i$ turned into bound variables), if the walk succeeds, \emph{dropped} if it fails by scope on a rigid path (it holds a
        forbidden eigenvariable that every instance keeps), and any other trouble suspends the whole equation.
        Then, if nothing was dropped and $G$ needs no lowering, $G\,\bar b$ stays as it is; otherwise, if some
        $b_j$ is functional, suspend; else prune and lower as in (b).
    \end{enumerate}
  \item \emph{Flexible $G\,\bar b$ in a soft context}: $G=F$ or $\ell(G)>\ell$ is trouble; otherwise the
    arguments are walked in a weak context and kept as they are. Inside an argument, the application of a
    variable bound inside that argument is treated as flexible too.
\end{enumerate}
Inside an argument of a non-pattern, case~(6c) reinterprets every failure: a scope failure of case~(1) in
a soft, non-weak context drops the argument, and any other $\bot$ or \emph{trouble} suspends the equation.
\end{definition}

\begin{tutorial}[Reading the contexts]
The context answers one question about the subterm being walked: \emph{will every instance of the right-hand
side still contain it?} At \texttt{top} and below rigid constructors (\texttt{rigid}, \texttt{firm}) the answer is
yes, because no substitution can remove a rigid constructor; so a forbidden eigenvariable there is a definite
failure, and a variable there can be lowered or pruned without losing solutions. Inside an argument of a
non-pattern $G\,\bar b$ the answer is \enquote{not necessarily}, since $G$ may become a function that ignores
that argument; so nothing is bound there and doubts turn into a suspension. \texttt{firm} and \texttt{rigid}
differ only in the occurs test (6a): below a constant, an eigenvariable outside $\bar a$ or a pair, $F$'s instance
cannot be a projection, so an inner $F\,\bar b$ with rigid arguments is fatal whatever their number; below a
bound variable or an $a_i$ at the head, the number matters. The next two tutorials give the reasons in detail.
\end{tutorial}

\begin{tutorial}[Strict and non-strict occurrences]
Pruning removes an argument from a variable; lowering makes a variable forget the eigenvariables it could see.
Both are sound only on \emph{strict} occurrences, those not inside an argument of a flexible term that is no
pattern (Dowek, Hardin, Kirchner and Pfenning \cite{DowekHardinKirchnerPfenning1996}). The reason is that
the flexible head $G$ may later turn out to drop that argument. Take
$\exists q\,\exists G\,\forall y\,\exists Y.\ q = G\,Y \wedge Y = y \wedge G = \lambda a.\,k$: it has the solution
$q=k$. If elimination of $q$ lowered $Y$ inside the argument of $G\,Y$, the later $Y=y$ would fail by scope
and the answer would be lost in that goal order, which breaks the commutativity of conjunction. So nothing is
bound inside such an argument: a variable to lower (case~(4)) or a flexible term to prune (case~(7)) is trouble
there, and trouble suspends the equation. The exception of case
(6c) is the one pruning that \emph{is} strict: an argument that holds, on a rigid path, an eigenvariable that $F$
may not see is held by every instance of that argument, so $G$ must drop it. This makes
$\exists q\,\exists F\,\forall x.\ q = F\,x\,x$ solvable as $F := \lambda a\,b.\,F'$, $q := F'$. Even that
pruning waits when an argument is, or may become, an abstraction: $G$ could apply it to another argument and
drop that one (\cref{sec:module-traces} shows $\exists q\,\forall x\,\exists G.\ q = G\,(\lambda w.\,k)\wedge G =
\lambda y.\,y\,x$, solved in both orders).
\end{tutorial}

\begin{tutorial}[When an occurrence of $F$ is fatal]
Case (6a) is Nipkow's size argument \cite{Nipkow1993}: $F\,\bar b$ inside the body of $F\,\bar a$ makes the
equation unsolvable only if no instance of $F\,\bar b$ can be smaller than the instance of $F\,\bar a$. That
holds when the $b_j$ are rigid (no substitution makes them reduce) and are as many as $\bar a$ plus the
binders on top ($F\,\bar a = \lambda\bar y.\,t'$ is $F\,\bar a\,\bar y = t'$ by η), as in
$F = \lambda x.\,g\,(F\,x)$ and \cref{ex:hounif-8i}. The count does not matter when the root of the body is a
constant, an eigenvariable not in $\bar a$ or a pair (context \texttt{firm}): no instance $\lambda\bar x.M$ of
$F$ can then be a projection, so $M$ has that root and $F\,\bar b$ sits strictly inside it. So the η-short
$F = \lambda x\,y.\,g\,(F\,y)$ fails as Miller and Nipkow fail it. Otherwise the instance may shrink and the
equation is outside $\Llam$: it is suspended. For instance $\forall x.\ F\,x = x\,(F\,(\lambda z.z))$ holds
with $F := \lambda a.\,a\,c$, and $\forall x\,\forall y.\ F\,x\,y = x\,(F\,y)$ holds with $F := \lambda w.w$.
\end{tutorial}

\subsection{The eager scope check}

\begin{definition}[Check]\label{def:module-check}
For a binding $X := t$ that micro has just made, $\textbf{Check}(X:=t)$ is $\mathrm{Elim}(\ell(X), -,
\langle\rangle, t)$ with no head to eliminate (\scm{μkanren-enforce}). If it answers $\bot$, \scm{=°} fails.
If it suspends, the binding is \emph{taken back}: \scm{μkanren-state-forget} unbinds $X$ and the equation
$X \doteq t$ is pushed on $E$. Otherwise its own bindings (lowerings, prunings) are kept.
\end{definition}
\textbf{Check} is NJK's binding rule with its lowering (\cref{def:rules-binding}): \enquote{$t$ contains no
constant with tag $>i$; change to $i$ the tags of the variables in $t$ whose tag is $>i$}
\cite[p.\,16]{NadathurJayaramanKwon1995}, extended from variables to flexible terms by pruning. Taking a binding
back is what Teyjus does with a binding whose check depends on a non-pattern. Its effect is an invariant:
\begin{proposition}[No captured eigenvariable]\label{prop:module-nocapture}
After every \scm{=°}, no binding $X := t$ in the substitution has an eigenvariable of level $>\ell(X)$ in the
value of $t$.
\end{proposition}
{\emergencystretch=2em
\noindent\emph{Argument.} Every binding made by micro goes through \textbf{Check} unless the fast paths of
\cref{sec:module-fast} prove it unnecessary; \textbf{Check} either fails, or rewrites the variables of $t$ so
that it holds, or unbinds $X$. The bindings that \textbf{Check} itself makes are lowerings and prunings,
which by construction respect the levels. The remaining gap is host code (\scm{project°} and \scm{set!}),
discussed in \cref{ch:roadmap}. \hfill\qedsymbol\par}

\subsection{Fast paths}\label{sec:module-fast}

\textbf{Check} on every binding would be costly on first-order search, so three cheap tests come first.
\begin{itemize}
  \item \textbf{No \scm{∀°} yet.} While no \scm{∀°} has run, the depth is 0 and $L$ is empty
    (\scm{μkanren-levels?} is false): every level is 0, no check can fail, none runs.
  \item \textbf{Plain terms.} \scm{μkanren-plain?} is an allocation-free walk, like micro's occurs check. It
    accepts $t$ when $t$ holds no eigenvariable or variable of a level above $\ell(X)$ and no application with a
    variable or abstraction head (which Elim might prune, lower or reduce). Only otherwise does Elim run.
  \item \textbf{Top-level variables.} A variable of level $\mathit{top}$, the highest level reached, can see
    every eigenvariable and every variable there is, so its bindings are skipped. The bindings made under
    \rulename{ξ} or \rulename{η} are checked anyway, by \scm{μkanren-scoped}.
\end{itemize}
The protocol's \texttt{atom?} hook belongs here as well: it tells micro's occurs check that an eigenvariable or
a bound variable holds no variable, so micro does not rebuild it as a slot list. On the type inference of the
Church numeral 400 these steps took the time from 463~ms to 330~ms (\scm{plain?}), 311~ms (the \texttt{top}
skip) and 254~ms (\texttt{atom?}).

\begin{pitfall}[The first \texttt{∀°} switches the check on for good]
The first fast path is about the whole branch, not the current scope: once a \scm{∀°} has run, levels exist
and every binding below \texttt{top} is checked, even after the \scm{∀°} has exited. First-order search that
follows a \scm{∀°} pays for it, measured at $+66\%$ on \scm{append°} backwards and $+127\%$ on open
\scm{append°}. Skipping the check at depth 0 would be unsound against host escapes; the fix planned is one walk
for micro's occurs check and the scope check together (\cref{ch:roadmap}).
\end{pitfall}

\subsection{Wake-up and termination}\label{sec:module-wake}

When \scm{=°} binds something and $E$ is not empty, \scm{μkanren-settle} retries every suspended item, oldest
first (\scm{μkanren-wake}). An item is retried by its kind:
\begin{itemize}
  \item an equation $u\doteq v$: if one side is now an unbound variable $X$ and the other no variable, $X$ is
    bound \emph{only once its check passes} (\scm{μkanren-retry/bind}); otherwise the equation is unified
    again, which may solve it, fail, or suspend it again;
  \item a type on a flexible term: decided if the head is now known, else suspended again;
  \item an absento: decided, kept, or dropped when levels entail it (\cref{sec:module-stores}).
\end{itemize}
The new bindings of a round are checked like those of \scm{=°}, and rounds repeat while a round binds. A
constraint still undecided is suspended again \emph{without} binding anything.

\begin{proposition}[Wake-up terminates]\label{prop:module-wake}
Suppose every hnf computed during the wake-up terminates. Then the loop of \scm{μkanren-settle} terminates.
\end{proposition}
\noindent\emph{Argument} (a sketch, not a mechanised proof). A round that binds nothing ends the loop, so it
suffices that only finitely many rounds bind. A round that binds does so for good: the only binding ever taken back is the one
that \textbf{Check} has just examined, and a retried $X \doteq t$ binds $X$ only after its check passed, so the
retry does not bind and unbind the same variable forever. The variables a round may create come from lowering
(case~(4)), whose new variable has a strictly smaller level than the one it replaces, and from pruning
(cases~(6b) and~(6c)), whose new variable either has a strictly smaller level (when $G$ is lowered, possibly with
\emph{extra} arguments) or the same level and strictly fewer arguments (\rulename{FF-same} likewise creates
$H$ at the level of $F$ with fewer arguments). So every new variable \emph{replaces} exactly one variable that
is bound for good, and is below it in the lexicographic order on (level, number of arguments), which is well
founded. A chain of replacements (a variable, the one that replaces it, the one that replaces that, \dots) is
therefore finite; with finitely many variables at the start, finitely many variables are ever created, each is
bound at most once, and only finitely many rounds can bind. This is the decreasing measure the design relies on.
One point a full proof would have to settle: the number of arguments is that of the occurrence being pruned,
which an untyped term does not fix once and for all. \hfill\qedsymbol

The hypothesis on hnf is not a formality: \scm{@} gives full β, and on untyped meta-terms full β may diverge
(\cref{ch:roadmap}). On pattern problems no such reduction happens (every contraction is a $\beta_0$ renaming),
and the rules are Miller's and Nipkow's algorithm with the prefix erased into levels: it terminates, fails only
when there is no unifier, and otherwise computes a most general one \cite[Thm~7.1, Prop.~7.3]{Miller1991},
\cite[Thm~2]{NadathurLinnell2005}. The suite checks the book's mgus verbatim: $U = \lambda x\,y.\,g\,(V'\,y)$ and
$V = \lambda y\,w.\,V'\,y$ for \cref{ex:hounif-8h}, $U = \lambda x.\,f\,(V'\,x)$ for \cref{ex:hounif-8k}, and
$G, H := \lambda w.\,T$ for $\lambda x\,\lambda y.\,G\,x = \lambda x\,\lambda y.\,H\,y$ (\cref{ex:hounif-ch4}).

\begin{pitfall}[Suspended is not solved]
Non-patterns are suspended, never enumerated: there is no Huet search \book{213}. An answer may therefore be
conditional on residual equations that are jointly unsatisfiable. $F\,a = b$ and $F\,a = c$ ($a$ a constant,
so no pattern) are both suspended, and the answer is printed with the two residuals; only a binding of $F$
reveals the conflict:
\begin{scheme}
(°->list #f (fresh° (q F) (=° q F) (=° (@ F 'a) 'b) (=° (@ F 'a) 'c)))
(°->list/datum (fresh° (q F) (=° q F) (=° (@ F 'a) 'b) (=° (@ F 'a) 'c)
                        (=° F (ƛ (z) 'b))))
\end{scheme}
\begin{answers}
((λ (α) (assert (≐ (@ α 'a) 'b)) (assert (≐ (@ α 'a) 'c)) α))
()
\end{answers}
This is floundering reported, not detected. The residuals are the honest output: they say exactly which
equations the search did not decide.
\end{pitfall}

\subsection{Worked traces}\label{sec:module-traces}

Each trace below poses a problem, follows the rules, and ends with the module's actual output. The probe used
here prints the depth, the levels of the given variables and $E$:
\begin{scheme}
(define ((probe label . vars) s)
  (display label) (display ": depth ") (display (μkanren-λ-depth s))
  (display ", levels ") (write (map (λ (v) (μkanren-λ-level v s)) vars))
  (display ", E ") (write (μkanren-λ-delayed s)) (newline)
  (✓° s))
\end{scheme}

\begin{example}[Warm-up: flexible–rigid, and \rulename{ξ}]\label{ex:module-warmup}
\emergencystretch=2em
Two equations that touch no pruning or lowering: $\exists q\,\forall x\,\forall y.\ q\,x\,y = g\,y\,x$, and
$\exists q.\ \lambda x.\,(g\ x) = \lambda y.\,(g\ (q\,y))$, where both $(g\ \dots)$ are Scheme lists.
\begin{scheme}
(°->list/datum (fresh° (q) (∀° (x y) (=° (@ q x y) `(g ,y ,x)))))
(°->list/datum (fresh° (q) (=° (ƛ (x) `(g ,x)) (ƛ (y) `(g ,(@ q y))))))
\end{scheme}
\begin{answers}
((ƛ (x y) (g y x)))
((ƛ (x) x))
\end{answers}
\emph{Trace of the first.} $\ell(q)=0$, $\ell(x)=1$, $\ell(y)=2$. The left side is flexible and $\langle x,y\rangle$
is a pattern for $q$; the right side is the list \scm{(g y x)}, rigid: \rulename{FR}. Elim walks the list in
context \texttt{firm}: \scm{g} is an atom and stays; $y$ is $a_2$ and becomes the bound variable of the second
binder; $x$ is $a_1$ and becomes that of the first (case~(1)). Nothing was bound on the way, and
$q := \lambda x\,y.\,(g\ y\ x)$, the binders named after the eigenvariables.

\emph{Trace of the second.} Both sides are abstractions: \rulename{ξ} opens them with one new eigenvariable $c$
of level $\infty$ (hint \texttt{x}), leaving \scm{(g c)} $\doteq$ \scm{(g (@ q c))}. These are lists, so micro
decomposes them and hands $c \doteq q\,c$ back to the hook. $\langle c\rangle$ is a pattern for $q$
($\infty > 0$), and Elim turns $c$ into the bound variable: $q := \lambda x.\,x$, a projection. The binding
was made under \rulename{ξ}, so \scm{μkanren-scoped} checks it at once; it holds no eigenvariable and passes.
\end{example}

\begin{example}[Pruning: problem 8.H \book{223--224}, \cref{ex:hounif-8h}]\label{ex:module-8h}
The problem (\cref{ex:hounif-8h}) is
\[
  \exists U\,\exists V\,\forall w\,\forall x\,\forall y.\ f\,(U\,x\,y) = f\,(g\,(V\,y\,w)).
\]
\begin{scheme}
(°->list/datum (fresh° (q u v) (=° q (list u v))
  (∀° (w x y) (probe "before" u v)
              (=° `(f ,(@ u x y)) `(f (g ,(@ v y w))))
              (probe "after" u v))))
\end{scheme}
\begin{answers}
before: depth 3, levels (0 0), E ()
after: depth 3, levels (0 0), E ()
(((ƛ (x y) (g (@ α y))) (ƛ (y w) (@ α y))))
\end{answers}
\emph{Trace.} $U$, $V$ have level 0 and $w,x,y$ levels 1, 2, 3. Both sides are lists, so micro decomposes
them: \scm{f} against \scm{f}, then $U\,x\,y \doteq g\,(V\,y\,w)$ with a λ-record on the left, handed to the
hook. The left side is a pattern ($x$, $y$ distinct, levels $2,3>0$) and the right is rigid: \rulename{FR},
$\mathrm{Elim}(0, U, \langle x,y\rangle, g\,(V\,y\,w))$. The root is the list \scm{(g ...)}: case (5),
context \texttt{firm}. Inside, $V\,y\,w$ is flexible with a pattern for $V$ (levels $3,1>0$): case~(6b). Keep $y$
(it is in $\bar a$, it becomes the bound variable of $y$); drop $w$ (not local, not in $\bar a$, and
$\ell(w)=1>0$); $\mathit{extra}$ is empty since $\ell(V)=0$ is not above $\ell=0$. Something was dropped, so
prune: $V := \lambda y\,w.\,V'\,y$, $\ell(V')=0$, and the body of $U$ is $g\,(V'\,y)$. Finally
$U := \lambda x\,y.\,g\,(V'\,y)$. Both are the book's mgu, with \texttt{α} for $V'$. The binders of $V$ are
named after the eigenvariables $y$, $w$ it was applied to.
\end{example}

\begin{example}[Flexible–flexible]\label{ex:module-ff}
Different heads, $\forall x\,\forall y.\ F\,x\,y = G\,y\,x$, and the same head, $\forall x\,\forall y.\ F\,x\,y =
F\,y\,x$:
\begin{scheme}
(°->list #f (fresh° (q f g) (=° q (list f g)) (∀° (x y) (=° (@ f x y) (@ g y x)))))
(°->list #f (fresh° (q) (∀° (x y) (=° (@ q x y) (@ q y x)))))
\end{scheme}
\begin{answers}
((λ (α) (cons (ƛ (x y) (@ α y x)) (cons α '()))))
((λ (α) (ƛ (x y) α)))
\end{answers}
\emph{Trace.} In the first, both sides are patterns with different heads: \rulename{FR} on $F$. In Elim, $G\,y\,x$
is case (6b) with both arguments in $\bar a$, kept; $\ell(G)=0$ is not above $\ell(F)=0$, so nothing is
pruned or lowered and $G$ stays unbound: $F := \lambda x\,y.\,G\,y\,x$. That is the book's mgu for different
heads \book{222} with the common variables listed as $y,x$: $F := \lambda x\,\lambda y.\,H\,y\,x$ and
$G := \lambda y\,\lambda x.\,H\,y\,x$, which is $H$ by η, up to the name of $H$. In the second, \rulename{FF-same}: the
positions where $x\,y$ and $y\,x$ agree are none, so $F := \lambda x\,y.\,H$ with $H$ new.
\end{example}

\begin{example}[A non-pattern, suspended and woken]\label{ex:module-nonpattern}
$\exists F\,\forall x.\ F\,(g\,x) = h\,(g\,x)$ is not in $\Llam$: the argument $g\,x$ is not an eigenvariable. It
has two solutions, $F := \lambda z.\,h\,z$ and $F := \lambda z.\,h\,(g\,x)$, but the second is illegal here
because $F$ is quantified outside $x$. The module neither guesses nor fails:
\begin{scheme}
(°->list #f (fresh° (q) (∀° (x) (=° (@ q `(g ,x)) `(h (g ,x))) (probe "inside" q))
                        (probe "after ∀" q)))
(°->list/datum (fresh° (q) (∀° (x) (=° (@ q `(g ,x)) `(h (g ,x))))
                           (probe "before F :=" q)
                           (=° q (ƛ (z) `(h ,z)))
                           (probe "after F :=" q)))
(°->list/datum (fresh° (q) (∀° (x) (=° (@ q `(g ,x)) `(h (g ,x))))
                           (=° q (ƛ (z) 'k))))
(°->list/datum (fresh° (q) (=° q (ƛ (z) `(h ,z)))
                           (∀° (x) (=° (@ q `(g ,x)) `(h (g ,x))))))
\end{scheme}
\begin{answers}
inside: depth 1, levels (0), E ((= (@ _0 (g (eigen x))) (h (g (eigen x)))))
after ∀: depth 0, levels (0), E ((= (ƛ (x) (@ _0 (g x))) (ƛ (x) (h (g x)))))
((λ (α)
    (assert
      (≐ (ƛ (x) (@ α (cons 'g (cons x '()))))
         (ƛ (x) (cons 'h (cons (cons 'g (cons x '())) '())))))
    α))
before F :=: depth 0, levels (0), E ((= (ƛ (x) (@ _0 (g x))) (ƛ (x) (h (g x)))))
after F :=: depth 0, levels (0), E ()
((ƛ (z) (h z)))
()
((ƛ (z) (h z)))
\end{answers}
\emph{Trace.} Inside the scope the left side is flexible but no pattern and the right is rigid:
\rulename{Delay}, and $E$ holds the equation, with the eigenvariable $x$ in it. At the exit of \scm{∀° (x)}
the equation mentions $x$, so it is \emph{discharged} (\cref{sec:module-discharge}): no variable in it has a
level $\ge\ell(x)=1$, so nothing is raised, and the equation is closed into
$\lambda x.\,F\,(g\,x) \doteq \lambda x.\,h\,(g\,x)$, equivalent by \rulename{ξ}. It is retried at once, opens
$x$ again as an $\infty$-level name, is still no pattern, and is suspended again. The run ends with it, and
the answer carries it as a residual \scm{(assert (≐ ...))}. In the second run the binding $F := \lambda z.\,h\,z$
wakes it: \rulename{ξ}, then β turns the left side into $h\,(g\,c)$, which equals the right side, and $E$ is
empty. In the third run $F := \lambda z.\,k$ makes the retry compare $k$ with $h\,(g\,c)$: a clash, so \scm{=°}
fails and there is no answer. In the fourth run the binding comes first and the equation is a plain β-redex,
solved with no suspension. The answer is the same in both goal orders.
\end{example}

\begin{example}[A binding taken back]\label{ex:module-forget}
$\exists q\,\exists G\,\exists H\,\forall y.\ q = G\,(H\,y)$. micro binds $q := G\,(H\,y)$ at once, as
\rulename{Var}. \textbf{Check} at level $\ell(q)=0$ then meets $G\,(H\,y)$, a flexible term whose argument
$H\,y$ is no eigenvariable: case (6c), walking $H\,y$ in a soft context. There $y$ (level 1, not allowed for a
level-0 variable) sits below the flexible head $H$, a weak context: trouble, not $\bot$, since $H$ may drop
$y$, and $G$ may drop all of $H\,y$. So the check suspends, the binding is taken back, and $q \doteq G\,(H\,y)$
waits in $E$:
\begin{scheme}
(°->list #f (fresh° (q G H) (∀° (y) (=° q (@ G (@ H y))) (probe "inside" q G H))))
(°->list/datum (fresh° (q G H) (∀° (y) (=° q (@ G (@ H y)))) (=° G (ƛ (a) 'k))))
\end{scheme}
\begin{answers}
inside: depth 1, levels (0 0 0), E ((= _0 (@ _1 (@ _2 (eigen y)))))
((λ (α β γ) (assert (≐ (ƛ (y) α) (ƛ (y) (@ β (@ γ y))))) α))
(k)
\end{answers}
At the exit of the scope the equation is closed over $y$ and reported, $q$ unbound, which is all one can say.
With $G := \lambda a.\,k$ the retry binds $q := k$, a plain term, and the residual is gone. The same happens in
every order of the goals.
\end{example}

\begin{example}[Occurs and scope failures, and a soft occurrence]\label{ex:module-occurs}\leavevmode
\begin{scheme}
(°->list #f (fresh° (q) (=° q `(f ,q))))                              ; X = f X
(°->list #f (fresh° (q) (∀° (w x y) (=° (@ q x y) `(g ,(@ q y w))))))  ; 8.I
(°->list #f (fresh° (q) (∀° (w x y) (=° (@ q x y) `(g ,w)))))          ; 8.J
(°->list #f (fresh° (F) (=° F (ƛ (x) `(g ,(@ F x))))))         ; F = λx.g (F x)
(°->list #f (fresh° (q F G) (=° q (list F G))
              (∀° (x) (=° (@ F x) (@ G (@ F x))))))
\end{scheme}
\begin{answers}
()
()
()
()
((λ (α β) (assert (≐ α (ƛ (x) (@ β (@ α x))))) (cons α (cons β '()))))
\end{answers}
The first is micro's occurs check, unchanged. In 8.I, $U\,y\,w$ occurs below the list \scm{(g ...)}, a
\texttt{firm} context, and its arguments are rigid: case (6a), $\bot$ (\cref{ex:hounif-8i}); it also has as
many arguments as $U\,x\,y$, so the count alone would fail it. In 8.J, $w$ has level~1, is not
an argument of $U$, and $U$ has level~0: case (1), $\bot$ by scope (\cref{ex:hounif-8j}). The fourth fails
micro's occurs check, the \texttt{occurs} hook passes it to Elim with no arguments, and $F\,x$ below \scm{g}
is case (6a) again. The last is the \emph{soft} occurs check: $F\,x$ occurs inside the argument of $G$, which may
be $\lambda z.z$, so the equation is suspended (and is closed over $x$ at the exit of the scope) instead of
failing.
\end{example}

\begin{example}[An argument that may be an abstraction]\label{ex:module-functional}
$\exists q\,\forall x\,\exists G.\ q = G\,(\lambda w.\,k) \wedge G = \lambda y.\,y\,x$ in both orders:
\begin{scheme}
(°->list/datum (fresh° (q) (∀° (x) (fresh° (G) (=° q (@ G (ƛ (w) 'k)))
                                               (=° G (ƛ (y) (@ y x)))))))
(°->list/datum (fresh° (q) (∀° (x) (fresh° (G) (=° G (ƛ (y) (@ y x)))
                                               (=° q (@ G (ƛ (w) 'k)))))))
\end{scheme}
\begin{answers}
(k)
(k)
\end{answers}
In the first order, \textbf{Check} of $q := G\,(\lambda w.k)$ at level 0 finds $G$ of level 1, which must be
lowered; but its argument is an abstraction that $G$ may apply to $x$ and so drop $G$'s own $x$. Lowering $G$
would not be strict, so case (6c) suspends. Then $G := \lambda y.\,y\,x$ wakes the equation, β gives
$(\lambda w.k)\,x = k$, and $q := k$. In the other order β happens at once. A version that lowered $G$ here
lost the answer in the first order.\sidenote{Third adversarial review.}
\end{example}

% ------------------------------------------------------------------------------------------------
\section{\texorpdfstring{\texttt{∀°}}{∀°} and \texorpdfstring{\texttt{⊃°}}{⊃°}}\label{sec:module-hh}

The book's goal-directed search for hereditary Harrop formulas consists of the rules \rulename{and}, \rulename{or}, \rulename{instan},
\rulename{generic}, \rulename{augment} and \rulename{backchain} \book{77}, \book{82}, \book{124} (\cref{fig:hh-rules}; the names are those of \cref{fig:rules-right}). Uniform proofs are
complete for intuitionistic provability of these formulas \cite{MillerNadathurPfenningScedrov1991}
(\cref{thm:hh-intuitionistic}), so implementing the rules faithfully \emph{is} the semantics. \rulename{and}, \rulename{or} and
$\exists$ are micro's \scm{and°}, \scm{or°} (interleaving) and \scm{fresh°}. This section gives the other
three as the module implements them.

\subsection{\texorpdfstring{\rulename{generic}}{GENERIC}}

\begin{mathpar}
\inferrule{\seq{\Sigma\cup\{c\}}{\Prog}{G[c/x]}}{\seq{\Sigma}{\Prog}{\forall x.\,G}}\;\rulename{generic}\quad(c\notin\Sigma)
\end{mathpar}
\scm{(∀° (x) g ...)} expands into \scm{(∀ª 'x (λ (x) (and° g ...)))}, and \scm{(∀° (x y) g ...)} into one
\scm{∀ª} per name, nested, so \scm{(∀° (w x y) ...)} gives $w,x,y$ the levels $d+1$, $d+2$, $d+3$. The Scheme
variable \scm{x} is bound to the eigenvariable, as \scm{ƛ} binds its placeholder. \scm{∀ª} takes a state $s$
whose λ record has depth $d$ and does the following.
\begin{enumerate}
  \item Record the level $d$ of every variable created so far (pad $L$ up to the variable count).
  \item Run the body \scm{(f c)} for a new eigenvariable $c$ of level $d+1$, in the state with depth $d+1$ and
    $\mathit{top} := \max(\mathit{top}, d+1)$.
  \item For each answer state $s_1$ of the body: record the level $d+1$ of the variables created inside (pad
    $L$ again), set the depth back to $d$, and \emph{discharge} the constraints on $c$
    (\cref{sec:module-discharge}); if discharge fails, drop the answer.
\end{enumerate}
\begin{scheme}
(define ((∀ª name f) s)
  (let* ((x  (μkanren-hh s))
         (d  (μkanren-hh-depth x))
         (x* (μkanren-hh-with (μkanren-hh-pad x (μkanren-state-vars-count s))
                              depth: (add1 d) top: (max (add1 d) (μkanren-hh-top x))))
         (c  (μkanren-eigen* name (add1 d))))
    (append-map§
      (λ (s1)
        (let1 (s2 (μkanren-with s1
                    (μkanren-hh-with (μkanren-hh-pad (μkanren-hh s1)
                                                     (μkanren-state-vars-count s1))
                                     depth: d)))
          (match/first (μkanren-discharge s2 c d)
            (#f (✗° s2))
            (,s3 (✓° s3)))))
      (δ ((f c) (μkanren-with s x*))))))
\end{scheme}
The shape is that of micro's \scm{andª} (\cref{ch:micro}): the body's stream is delayed with \scm{δ}, and
\scm{append-map§} post-processes it answer by answer, so the restoring step runs on every answer, however
late the stream produces it. Steps 1 and 3 pad $L$ because a variable's level is \emph{implicit} (the current
depth) until $L$ is extended past it: padding fixes it before the depth changes.

The signature $\Sigma$ is not carried. What carries the quantifier structure is the level of each variable:
a variable created at depth $d$ stands for an $\exists$ after the first $d$ universals of its branch (Miller's
mixed prefix \cite{Miller1991}, linearised by creation), and \textbf{Check} forbids exactly the substitutions
that would violate that prefix. NJK prove this tagging sound and complete (\cref{prop:rules-ui},
\cite[Prop.\,4.1]{NadathurJayaramanKwon1995}).

\begin{tutorial}[Why the depth must be restored on every answer]
A creation timestamp would tell a variable created before $c$ from one created after it, and one might think
that suffices: older variables may not see $c$. But a variable created \emph{after} the \scm{∀°} has exited
may not see $c$ either. With the depth restored on every answer, such a variable gets level $d < \ell(c)$, so
it cannot take $c$. This is λProlog's universe register \cite{NadathurJayaramanKwon1995}. The two classic
cases, \book{85}:
\begin{scheme}
(list (°->list #f (fresh° (q) (∀° (x) (=° q x))))         ; ∃Y∀x. Y = x
      (°->list #f (fresh° (q) (∀° (x) (fresh° (y) (=° y x)))
                              (=° q 'ok))))                ; ∀x∃Y. Y = x
\end{scheme}
\begin{answers}
(() ((λ () 'ok)))
\end{answers}
In the first, $q$ has level~0 and $x$
level~1: \textbf{Check} fails by case (1). In the second, $y$ is created inside, at level~1, and may take $x$.
\end{tutorial}

\begin{example}[Lowering, seen from the state]\label{ex:module-lower}
A variable created inside the scope and placed inside the value of an outer one must forget the scope:
\begin{scheme}
(define ((probe . vars) s)
  (for-each (λ (v) (let1 (w (μkanren-state-find v s))
                     (write v) (display " ↦ ") (write w)
                     (when (μkanren-var? w)
                       (display " level ") (display (μkanren-λ-level w s)))
                     (display "; ")))
            vars)
  (newline) (✓° s))
(°->list #f (fresh° (A) (∀° (x) (fresh° (B) (probe A B)
                                             (=° A `(arr ,B c))
                                             (probe A B)))))
\end{scheme}
\begin{answers}
_0 ↦ _0 level 0; _1 ↦ _1 level 1;
_0 ↦ (arr _1 c); _1 ↦ _2 level 0;
((λ (α) (cons 'arr (cons α (cons 'c '())))))
\end{answers}
$A$ (\texttt{\_0}) has level~0 and $B$ (\texttt{\_1}) level~1. Binding $A := (\mathtt{arr}\ B\ c)$ triggers
\textbf{Check} at level~0, which meets the unbound $B$ of level~$1>0$: case (4), lower. A new variable
\texttt{\_2} of level~0 is created and $B := \mathtt{\_2}$; \textbf{Check} eliminates no head and has no
arguments $\bar a$ to raise over, so the raising part is empty. This is NJK's \enquote{change to $i$ the tags of the variables in $t$ whose tag is $>i$}, done by a
binding instead of a mutation, so that it is undone by backtracking for free.
\end{example}

\subsection{Discharge at the exit of a scope}\label{sec:module-discharge}

When \scm{∀°} exits, some constraints may still mention $c$: suspended equations, generic disequalities of N,
absentos of micro's T. A residual that mentions $c$ outside its scope would be meaningless, and the printer
must never show an eigenvariable free. \scm{μkanren-discharge} closes them, back at depth $d$:
\begin{enumerate}
  \item \textbf{Raise.} Every unbound variable $Y$ of the equations and disequalities to be closed with
    $\ell(Y)\ge\ell(c)$ is raised:
    $Y := Y'\,c$ with $\ell(Y')=d$. Inside the scope $Y$ could depend on $c$; $Y'$ lives outside it, so the
    dependency becomes an explicit argument: $\forall c\,\exists Y$ is turned into $\exists Y'\,\forall c$
    with $Y = Y'\,c$. This is Miller's raising \book{107} (\cref{prop:lambda-raising}), a bijection
    on solutions.
  \item \textbf{Close.} A suspended equation $u \doteq v$ that mentions $c$ stands for
    $\forall c\,(u = v)$ and becomes $\lambda c.u \doteq \lambda c.v$, equivalent by \rulename{ξ}. A generic
    disequality is closed the same way, under \rulename{generic}'s reading of $c$ as a new name. An absento of T whose
    \emph{term} mentions $c$ is closed too, since absence in $\lambda c.u$ is absence in $u$. An absento whose
    \emph{datum} mentions $c$ moves to $E$ as \scm{(absent d u)}, unless the levels entail it: if $u$ does not
    hold $c$ and all its variables have levels below $\ell(c)$, no instance of $u$ can ever hold $c$, and the
    constraint is dropped.
  \item \textbf{Retry.} The raise bindings go through \scm{μkanren-post-unify}, which wakes $E$, so the closed
    constraints are retried. With no raise, the closed equations are retried anyway.
  \item \textbf{Repeat.} A raise may bring $c$ into more constraints (an absento of T on a raised $Y$ now
    watches $Y'\,c$; a prefix of D on $Y$ moves to N holding $c$). Discharge repeats until no raise is left.
\end{enumerate}
The constraints that cannot be closed by an abstraction (a type on $F\,x$; an absento whose datum is $c$ on a
term that may still hold $c$) keep $c$ and print under a \scm{∀} (\cref{sec:module-reify}).

\begin{example}[Discharge: raising, closing, retrying]\label{ex:module-discharge}\leavevmode
\begin{scheme}
(define ((probe label . vars) s)   ; prints values and the levels of unbound variables
  (display label) (display ": depth ") (display (μkanren-λ-depth s))
  (for-each (λ (v) (display ", ") (write v) (display " ↦ ")
                   (write (μkanren-state-find/value v s))
                   (let1 (w (μkanren-state-find v s))
                     (when (μkanren-var? w)
                       (display " level ") (display (μkanren-λ-level w s)))))
            vars)
  (display ", E ") (write (μkanren-λ-delayed s)) (newline)
  (✓° s))
(°->list #f (fresh° (q) (∀° (a) (fresh° (f) (=° (@ f a) `(g ,a ,a))
                                             (probe "inside" f)))
                        (=° q 'ok)))
(°->list #f (fresh° (q) (∀° (x) (=° (@ q x) `(g ,q)) (probe "inside" q))))
\end{scheme}
\begin{answers}
inside: depth 1, _1 ↦ _1 level 1, E ((= (@ _1 (eigen a)) (g (eigen a) (eigen a))))
((λ (α)
    (assert (≐ (ƛ (a) (@ α a a)) (ƛ (a) (cons 'g (cons a (cons a '()))))))
    'ok))
inside: depth 1, _0 ↦ _0 level 0, E ((= (@ _0 (eigen x)) (g _0)))
()
\end{answers}
In the first, $f$ is created inside the scope of $a$ (level 1), so $f\,a$ is no pattern (the book's 8.A, four
unifiers and no mgu, \book{213}) and is suspended. At the exit, $f$ has level $1\ge\ell(a)$: it is raised,
$f := f'\,a$ with $\ell(f')=0$, and the equation closes into $\lambda a.\,f'\,a\,a \doteq \lambda a.\,g\,a\,a$. That is
still no pattern (the argument $a$ is repeated), so it is reported. In the second, $F\,x \doteq g\,F$ is a pattern
equation, but $F$ itself occurs with no argument in its own body, case (4) with $\bar a \neq \langle\rangle$:
suspend. At the exit nothing needs raising ($\ell(F)=0$), the equation closes into
$\lambda x.\,F\,x \doteq \lambda x.\,g\,F$, which by η is $F \doteq \lambda x.\,g\,F$, and the retry fails the occurs
check: no answer, as it must be.
\end{example}

\subsection{\texorpdfstring{\rulename{augment}}{AUGMENT}}

\begin{mathpar}
\inferrule{\seq{\Sigma}{\Prog\cup\{D\}}{G}}{\seq{\Sigma}{\Prog}{D\imp G}}\;\rulename{augment}
\end{mathpar}
\scm{(⊃° (D ...) g ...)} expands into \scm{⊃ª}, which pushes the clauses on Γ for the body and gives the old Γ
back to every answer of the body:
\begin{scheme}
(define ((⊃ª clauses g) s)
  (let* ((x (μkanren-hh s))
         (Γ (μkanren-hh-Γ x))
         (Γ* (fold-right
               (λ (c Γ*)
                 (let1 (old (cond ((assq (car c) Γ*) => cdr) (else '())))
                   (cons (cons (car c) (cons (cdr c) old))
                         (if (null? old) Γ* (alist-delete (car c) Γ* eq?)))))
               Γ clauses)))
    (append-map§
      (λ (s1) (✓° (μkanren-with s1 (μkanren-hh-with (μkanren-hh s1) Γ: Γ))))
      (δ (g (μkanren-with s (μkanren-hh-with x Γ: Γ*)))))))
\end{scheme}
A clause is a pair \scm{(relation . procedure)}: the relation is the dynamic relation's own procedure, used
as an \scm{eq?} key, and the procedure maps the list of a call's arguments to the goal that the clause makes of
them. Since Γ lives in the state, \scm{or°} branches cannot interfere, and restoring it on every answer is the
stack discipline of the book \book{78}, \book{83}. Assuming a clause for a relation already in Γ replaces that
relation's entry by one with the new clause in front of the old ones, so a lookup costs the number of distinct
assumed relations, not the number of clauses. The \scm{fold-right} puts the first clause of a list on top.

\subsection{\texorpdfstring{\rulename{backchain}}{BACKCHAIN} and \texorpdfstring{\texttt{define-dynamic-relation}}{define-dynamic-relation}}

In λProlog every predicate can be extended by an assumption. Here only relations made by
\scm{define-dynamic-relation} can:
\begin{scheme}
(define-syntax define-dynamic-relation
  (syntax-rules ()
    ((_ (name arg ...))
     (define name
       (μkanren-dynamic
         (λ (arg ...) (λ (s) (δ ((μkanren-hyps° name (list arg ...) #f) s)))))))
    ((_ (name arg ...) g ...)
     (define name
       (μkanren-dynamic
         (λ (arg ...)
           (λ (s) (δ ((μkanren-hyps° name (list arg ...) (and° g ...)) s)))))))))

(define ((μkanren-hyps° key args static) s)
  (let* ((gs (map (λ (c) (c args))
                  (cond ((assq key (μkanren-hh-Γ (μkanren-hh s))) => cdr) (else '()))))
         (gs (if static (append gs (list static)) gs)))
    (if (null? gs) '() ((reduce-right orª #f gs) s))))
\end{scheme}
A call disjoins, in one flat \scm{or°}, the assumed clauses for the relation, \emph{most recent first}, and then
the static body. If $C_1,\dots,C_m$ are the clauses assumed for $r$ in the current Γ, in the order they were
assumed, and $B_r$ is the static body of $r$ (absent when it has none), then
\[
  r\,\bar t \;\rightsquigarrow\; C_m(\bar t) \;\vee\; \dots \;\vee\; C_1(\bar t) \;\vee\; B_r(\bar t),
\]
where $C_i(\bar t)$ unifies $\bar t$ with the head of $C_i$, its own \scm{∀} variables renamed and its free
variables shared (see the clause syntax below), and runs its body. This is
\rulename{backchain} (\cref{fig:hh-rules}) with the choice of clause made by \scm{or°} instead of by a selection rule; a
relation with neither hypotheses nor static body fails. When every clause answers at once, the answers come in λProlog's order; otherwise
interleaving decides. The relation is tagged with CHICKEN's \scm{extend-procedure}, so \scm{μkanren-dynamic?}
is an $O(1)$ test and no global registry is needed. Assuming a clause for a relation that is not dynamic is an
error, not a silent no-op:
\begin{scheme}
(define-dynamic-relation (n° x) (=° x 1))
(list (°->list/datum (fresh° (q) (⊃° ((n° 2)) (⊃° ((n° 3)) (n° q)))))
      (°->list/datum (fresh° (q) (⊃° ((n° 2) (n° 3)) (n° q)))))
(define (error-message thunk)     ; the message of the error `thunk` raises
  (handle-exceptions exn ((condition-property-accessor 'exn 'message #f) exn)
    (begin (thunk) #f)))
(define-relation (r° x) (=° x 1))
(error-message (λ () (°->list #f (fresh° (q) (⊃° ((r° 2)) (r° q))))))
\end{scheme}
\begin{answers}
((3 2 1) (2 3 1))
"not a dynamic relation, see define-dynamic-relation"
\end{answers}
The two orders are the book's \book{78}: \lp{p 2 => p 3 => p X} gives 3, 2, 1, and
\lp{(p 2 & p 3) => p X} gives 2, 3, 1. In the first query the inner \scm{⊃°} runs last, so Γ maps \scm{n°} to
the clauses for 3 and 2, in this order, and the call disjoins $q=3$, $q=2$ and the static $q=1$. In the second,
\scm{fold-right} puts the first clause of the list on top, so the clause for 2 comes before the one for 3.

\subsection{Clause syntax}

\begin{longtable}{@{}p{0.36\fullwidthlen}p{0.58\fullwidthlen}@{}}
\toprule
clause $D$ & meaning and expansion \\
\midrule
\endhead
\scm{(rel t ...)} & a fact; \scm{(λ (args) (fresh° (v ...) (=° args (list t ...))))} \\
\scm{((rel t ...) ← g ...)} & a rule, infix \\
\scm{(← (rel t ...) g ...)} & the same rule, prefix; rewritten to the infix form \\
\scm{(∀ (w ...) D)} & $D$ with its own variables \scm{w ...}, accumulated into the \scm{fresh°} of the expansion \\
\bottomrule
\end{longtable}
\noindent A rule expands to \scm{(μkanren-clause* rel (λ (args) (fresh° (v ...) (=° args (list t ...)) g ...)))},
and \scm{μkanren-clause*} raises the error above when \scm{rel} is not dynamic. Two kinds of variables appear
in a clause, as in the book \book{86--88}:
\begin{itemize}
  \item its \emph{free} logic variables are those of the closure: shared with the context, and among all uses of
    the clause;
  \item its \emph{own} \scm{∀} variables are made by the \scm{fresh°} inside the procedure, so they are renamed at
    each use, at the current level. That is NJK's \enquote{new tagged instance} of a clause
    \cite[p.\,14]{NadathurJayaramanKwon1995}.
\end{itemize}
\begin{scheme}
(define-dynamic-relation (h° x))
(list (°->list/datum (fresh° (q y) (⊃° ((h° `(,y ,y))) (h° '(1 1)) (h° '(2 2)))))
      (°->list/datum (fresh° (q) (⊃° ((∀ (z) (h° `(,z ,z)))) (h° '(1 1)) (h° '(2 2)))
                                 (=° q 'ok))))
\end{scheme}
\begin{answers}
(() (ok))
\end{answers}
With the shared $y$, the first use binds $y:=1$ and the second use cannot match \scm{(2 2)}; with the clause's
own $z$, each use gets a new $z$.

\begin{example}[Levels reach into clauses {\cite[p.\,16]{NadathurJayaramanKwon1995}}]\label{ex:module-njk}
\emergencystretch=2em
With the program $\forall z\,(q\,z \imp p\,(d\,z))$, the goal $\exists x\,\forall y\,(q\,y \imp p\,x)$ must fail
and $\forall y\,\exists x\,(q\,y \imp p\,x)$ must succeed (\cref{ex:rules-scope}):
\begin{scheme}
(define-dynamic-relation (q° z))
(define-dynamic-relation (p° x) (fresh° (z) (=° x `(d ,z)) (q° z)))
(list (°->list #f (fresh° (x) (∀° (y) (⊃° ((q° y)) (p° x)))))
      (°->list #f (fresh° (r) (∀° (y) (fresh° (x) (⊃° ((q° y)) (p° x)))) (=° r 'ok))))
\end{scheme}
\begin{answers}
(() ((λ () 'ok)))
\end{answers}
In the first goal, $x$ has level~0. Backchaining on \scm{p°} creates the clause variable $z$ inside the scope of
$y$ (level~1) and binds $x := (d\ z)$; \textbf{Check} lowers $z$ to level~0. The hypothesis \scm{(q° y)} then
asks $z = y$, a level-0 variable against a level-1 eigenvariable: scope failure. Without lowering, $z$ would
keep level~1 and take $y$, and $x = (d\ y)$ would leak $y$ out of its scope.
\end{example}

\subsection{Worked example: type inference}\label{sec:module-typeof}

The book's type inference for untyped λ-terms (\cref{sec:lambdatree-typeof}, \book{190--191}) uses all three
rules at once:
\begin{lprolog}
typeof (app M N) A :- typeof M (arr B A), typeof N B.
typeof (abs M) (arr A B) :- pi x\ typeof x A => typeof (M x) B.
\end{lprolog}
In the module's syntax ({\small\texttt{test/programs.scm}} of {\small\texttt{doc/hol-kanren/prototype/}}):
\begin{scheme}
(define-dynamic-relation (typeof° m t)
  (cond°
    ((fresh° (M N B) (=° m `(app ,M ,N)) (typeof° M `(arr ,B ,t)) (typeof° N B)))
    ((fresh° (R A B) (=° m `(abs ,R)) (=° t `(arr ,A ,B))
       (∀° (x) (⊃° ((typeof° x A)) (typeof° (@ R x) B)))))))
\end{scheme}
The object term \lp{abs x\ M} is \scm{`(abs ,(ƛ (x) M))}; \lp{(M x)} is \scm{(@ R x)}, a β-redex once $R$ is
bound. The answers for S, I and $\lambda x.\,x\,x$ (from \texttt{examples.out} in the same directory):
\begin{answers}
(((arr (arr α (arr β γ)) (arr (arr α β) (arr α γ)))) ((arr α α)) ())
\end{answers}
To see the rules at work, here is the same program with a trace goal that prints the depth and the walked
arguments, with the level of each unbound variable after \texttt{@}. The trace in the static body runs only
when the static clauses are tried, after the hypotheses:
\begin{scheme}
(define ((trace label . ts) s)
  (display label) (display " [depth ") (display (μkanren-λ-depth s)) (display "]")
  (for-each (λ (t) (display " ") (write (μkanren-state-find/value t s))
                   (let1 (w (μkanren-state-find t s))
                     (when (μkanren-var? w)
                       (display "@") (display (μkanren-λ-level w s)))))
            ts)
  (newline) (✓° s))
(define-dynamic-relation (typeof° m t)
  (trace "static typeof°" m t)
  (cond°
    ((fresh° (M N B) (=° m `(app ,M ,N)) (trace " app" M N B)
       (typeof° M `(arr ,B ,t)) (typeof° N B)))
    ((fresh° (R A B) (=° m `(abs ,R)) (=° t `(arr ,A ,B)) (trace " abs" R A B)
       (∀° (x) (⊃° ((typeof° x A)) (typeof° (@ R x) B)))))))
(°->list/datum (fresh° (q) (typeof° `(abs ,(ƛ (x) `(abs ,(ƛ (y) x)))) q)))
(°->list/datum (fresh° (q) (typeof° `(abs ,(ƛ (x) `(app ,x ,x))) q)))
\end{scheme}
\begin{answers}
static typeof° [depth 0] (abs (ƛ (x) (abs (ƛ (y) x)))) _0@0
 abs [depth 0] (ƛ (x) (abs (ƛ (y) x))) _2@0 _3@0
static typeof° [depth 1] (abs (ƛ (y) (eigen x))) _3@0
 abs [depth 1] (ƛ (y) (eigen x)) _7@0 _8@0
static typeof° [depth 2] (eigen x) _8@0
((arr α (arr β α)))
static typeof° [depth 0] (abs (ƛ (x) (app x x))) _0@0
 abs [depth 0] (ƛ (x) (app x x)) _2@0 _3@0
static typeof° [depth 1] (app (eigen x) (eigen x)) _3@0
 app [depth 1] (eigen x) (eigen x) _6@1
static typeof° [depth 1] (eigen x) (arr _6 _3)
static typeof° [depth 1] (eigen x) _7@0
()
\end{answers}

\begin{tutorial}[$\lambda x.\lambda y.\,x$, step by step]
\begin{enumerate}
  \item \emph{Depth 0, Γ empty.} The query $T$ is \texttt{\_0}. No hypothesis for \scm{typeof°}, so the static
    body runs. The \texttt{abs} clause creates $R,A,B$ (\texttt{\_1}, \texttt{\_2}, \texttt{\_3}, level~0),
    binds $R := \lambda x.\,\mathtt{abs}\,(\lambda y.\,x)$ and $T := (\mathtt{arr}\ A\ B)$.
  \item \rulename{generic}. \scm{∀° (x)} records levels, raises the depth to~1 and makes the eigenvariable $x$ of
    level~1.
  \item \rulename{augment}. \scm{⊃°} pushes $\langle\texttt{typeof°}, \lambda\mathit{args}.\ \mathit{args} \doteq (x\ A)\rangle$ on Γ.
  \item \emph{β.} The goal is \scm{(typeof° (@ R x) B)}; \scm{@} reduces the redex to
    $\mathtt{abs}\,(\lambda y.\,x)$, with $x$ now the eigenvariable, printed \scm{(eigen x)}.
  \item \rulename{backchain}, depth 1. The hypothesis is tried first: $(\mathtt{abs}\ \dots) \doteq x$ clashes.
    Then the static body: the \texttt{abs} clause creates $R',A',B'$ at level~1 and binds
    $B := (\mathtt{arr}\ A'\ B')$. Here \textbf{Check} runs: $B$ has level~0, $A'$ and $B'$ level~1, so both are
    \emph{lowered}, to new variables \texttt{\_7}, \texttt{\_8} of level~0 (the trace shows
    \texttt{\_7@0 \_8@0}, the walked values). This is right: the type of the inner abstraction is part of the
    outer type, which lives outside the scope of $x$.
  \item \rulename{generic}, \rulename{augment} again. Depth~2, eigenvariable $y$ of level~2, Γ now holds two clauses for
    \scm{typeof°}, $(y\ A')$ on top of $(x\ A)$.
  \item \rulename{backchain}, depth 2. The goal is \scm{(typeof° x B')}. The most recent hypothesis gives
    $x \doteq y$, a clash of two eigenvariables. The older one gives $B' \doteq A$: a level-0 variable bound to
    another, fine. The static body is tried as well (the last trace line) and fails, since $x$ is neither an
    application nor an abstraction.
  \item \emph{Exits.} Each \scm{∀°} gives its answers back their depth and discharges nothing (no constraint
    mentions $x$ or $y$); each \scm{⊃°} gives Γ back. The answer is
    $T = \mathtt{arr}\ A\ (\mathtt{arr}\ A'\ A)$, printed \scm{(arr α (arr β α))}.
\end{enumerate}
\end{tutorial}

\begin{tutorial}[$\lambda x.\,x\,x$ fails by the occurs check]
At depth~1, with the hypothesis $\mathtt{typeof}\ x\ A$, the goal is
$\mathtt{typeof}\ (\mathtt{app}\ x\ x)\ B$. The \texttt{app} clause binds $M,N := x$ and creates its $B''$
(\texttt{\_6}) at level~1. The first subgoal $\mathtt{typeof}\ x\ (\mathtt{arr}\ B''\ B)$ uses the hypothesis:
$A := (\mathtt{arr}\ B''\ B)$, and since $\ell(A)=0$, \textbf{Check} lowers $B''$ to the level-0
\texttt{\_7}. The second subgoal $\mathtt{typeof}\ x\ B''$ (printed with $B''$ walked to \texttt{\_7}) asks
$\mathtt{\_7} \doteq A = (\mathtt{arr}\ \mathtt{\_7}\ B)$: micro's occurs check fails, and since the term holds
no λ-term, the \texttt{occurs} hook does not intervene. The static attempts fail too, and there is no answer,
the book's \enquote{no} \book{191}.
\end{tutorial}

% ------------------------------------------------------------------------------------------------
\section{micro's stores on λ-terms}\label{sec:module-stores}

λProlog has no disequality, no absento and no type constraints such as \scm{symbol°} (its types are
static). \micro{} has all three (\cref{sec:micro-stores}), and
every relation built on them must keep working when its arguments are λ-terms. The readings adopted here extend
micro's semantics over λProlog's term model: equality is $\alpha\beta\eta$-equality, a flexible term is
undecided until its head is known, and an eigenvariable is, by default, a new name.

\subsection{Disequality: the \texorpdfstring{\texttt{diseq}}{diseq} hook and the store N}

\begin{definition}[Trial]\label{def:module-trial}
The module \emph{claims} a disequality $u \neq v$ when either side, walked, holds a λ-record or an
eigenvariable, and also, under the extensional reading once a \scm{∀°} has run, a first-order one (in D an
absento of T could absorb it, and absence reads eigenvariables generically). Its \emph{trial} unifies $u \doteq v$ in the current state, with \textbf{Check}, and classifies
the result (\scm{μkanren-trial}):
\begin{description}[style=nextline,leftmargin=2.2em]
  \item[\texttt{unsat}] no unifier: the disequality holds for good and is dropped;
  \item[\texttt{entailed}] a unifier that binds no variable older than the trial (or binds older variables only
    to distinct variables the trial itself created) and suspends nothing: $u=v$ already holds, and \scm{≠°}
    fails;
  \item[kept] otherwise: the pair is stored in N as \scm{(reading u . v)} and re-tried after every \scm{=°}.
\end{description}
\end{definition}
The variables a trial creates (by pruning, by \rulename{FF-same}, by lowering) are \emph{existential inside the
negation}: $\lambda x\,y.\,F\,x \neq \lambda x\,y.\,G\,y$ means $\neg\exists H.\ F = \lambda x.H \wedge G =
\lambda y.H$. A prefix of micro's D would wrongly store them as free variables, which is one reason N is a
separate store.
\begin{scheme}
(°->list #f (fresh° (q) (=° q (ƛ (x) x)) (≠° q (ƛ (y) y))))     ; entailed mod α
(°->list #f (fresh° (q) (≠° (@ (ƛ (x) `(g ,x)) 'a) '(g a))))    ; entailed mod β
(°->list #f (fresh° (q) (≠° (ƛ (x) x) (ƛ (x) 'k)) (=° q 'ok)))  ; holds for good
(°->list/datum (fresh° (q F) (=° q F) (≠° (@ F 'a) 'a)
                        (=° F (ƛ (z) z))))                 ; woken, violated
(°->list/datum (fresh° (q F) (=° q F) (≠° (@ F 'a) 'a)
                        (=° F (ƛ (z) 'b))))                ; woken, holds
(°->list #f (fresh° (q f g) (=° q (list f g))
              (≠° (ƛ (x y) (@ f x)) (ƛ (x y) (@ g y)))))
\end{scheme}
\begin{answers}
()
()
((λ () 'ok))
()
((ƛ (z) b))
((λ (α β)
    (begin (deny (∃ (_3) (and (equal? α (ƛ (x) _3)) (equal? β (ƛ (y) _3))))))
    (cons α (cons β '()))))
\end{answers}
Line by line: the first trial succeeds under \rulename{ξ} binding nothing, so it is \texttt{entailed} and
\scm{≠°} fails; the second never reaches the module, since \scm{@} has already β-reduced its left side to the
list \scm{(g a)} and micro's own \scm{≠°} fails on two equal terms; the third trial clashes ($x$ against $k$ under \rulename{ξ}), so the disequality is \texttt{unsat} and
dropped. In the fourth and fifth the trial $F\,a \doteq a$ is no pattern ($a$ is a constant) and suspends, so the
pair is kept in N; the binding of $F$ retries it, and it is violated by $\lambda z.z$ and holds for good
with $\lambda z.b$. In the last, the trial is the flexible–flexible problem of \cref{ex:module-ff} under two
binders: it binds the older $F$ and $G$ to $\lambda x.H$ and $\lambda y.H$ with a new $H$, so the pair is kept,
and the answer denies exactly that mgu, with $H$ (\texttt{\_3}) existential inside the \scm{deny}.

\paragraph{Migration from D to N.} A prefix of D born first-order may come to hold a λ-term or an eigenvariable
later, at any depth below its variables. The \texttt{post} hook moves it to N, as the tuple disequality it
stands for, \emph{before} micro re-verifies D. So D never holds a λ-term, and micro's plain trial never decides a
λ-equation without the scope check and the reading. The probe prints N, D and E:
\begin{scheme}
(define ((probe label) s)
  (display label) (display ": N ") (write (μkanren-λ-disequalities s))
  (display ", D ") (write (μkanren-state-D s))
  (display ", E ") (write (μkanren-λ-delayed s))
  (newline) (✓° s))
(°->list #f (fresh° (q Y) (≠° q (list Y)) (probe "before")
                          (=° Y (ƛ (z) z)) (probe "after")))
\end{scheme}
\begin{answers}
before: N (), D (((_0 _1))), E ()
after: N ((generic (_0) (_1))), D (), E ()
((λ (α) (begin (deny (equal? α (cons (ƛ (z) z) '())))) α))
\end{answers}
To keep first-order search cheap, the hook looks at no more than 64 nodes of the new bindings to decide whether
they may hold a λ-term; past that bound it walks D and A instead.\sidenote{\Cref{ch:roadmap} has the measurements.}

\paragraph{Subsumption by A.} A trial that binds a variable typed in A to a rigid λ-term or an eigenvariable
of which the type's predicate is false has no solution; \scm{symbol°} and \scm{number°} hold of neither. So the disequality is dropped, as micro
drops \scm{(≠° q 5)} under \scm{(symbol° q)}:
\begin{scheme}
(°->list #f (fresh° (q) (symbol° q) (≠° q (ƛ (x) x))))
\end{scheme}
\begin{answers}
((λ (α) (assert (every (μ v (symbol? v)) (list α))) α))
\end{answers}

\paragraph{Two readings of $\forall$.} What does $\forall x.\ q \neq x$ mean? λProlog has no $\neq$, but it has
\lp{pi x\ not (q = x)}, which succeeds: $x$ is a new name, distinct from everything. The parameter
\scm{μkanren-λ-reading} chooses, and each disequality keeps the reading in force when it was posted.
\begin{itemize}
  \item \scm{'generic} (the default) reads an eigenvariable as \rulename{generic}'s new name, ∇-like, as αKanren's noms
    \cite{ByrdFriedman2007}. At the exit of the scope the constraint is closed: $\forall x.\ F\,x \neq x$
    becomes $F \neq \lambda x.x$, consistently whichever comes first, the binding or the $\neq$.
  \item \scm{'extensional} reads an eigenvariable as any term of its universe, as in the miniKanren work on
    universal quantification \cite{Moiseenko2019,JinEtAl2021}. The trial replaces each eigenvariable by a new
    variable of its level, after raising the variables that may depend on it, and counts as entailed when it
    binds older variables only to distinct such variables. Then $\forall x.\ q \neq x$ fails, since $x$ could
    be $q$.
\end{itemize}
\begin{scheme}
(°->list #f (fresh° (q) (∀° (x) (≠° q x))))
(°->list #f (fresh° (q F) (=° q F) (∀° (x) (≠° (@ F x) x))))
(°->list #f (fresh° (q F) (=° q F) (∀° (x) (≠° (@ F x) x)) (=° F (ƛ (z) z))))
(°->list #f (fresh° (q F) (=° q F) (∀° (x) (≠° (@ F x) x)) (=° F (ƛ (z) 'c))))
(parameterize ((μkanren-λ-reading 'extensional))
  (list (°->list #f (fresh° (q) (∀° (x) (≠° q x))))
        (°->list #f (fresh° (q F) (=° q F) (∀° (x) (≠° (@ F x) x))))
        (°->list #f (fresh° (q F) (=° q F) (∀° (x) (≠° (@ F x) x))
                                (=° F (ƛ (z) 'c))))))
\end{scheme}
\begin{answers}
((λ (α) α))
((λ (α) (begin (deny (equal? α (ƛ (x) x)))) α))
()
((λ () (ƛ (z) 'c)))
(() ((λ (α) (∀ (x) (begin (deny (≐ (@ α x) x)))) α)) ())
\end{answers}
Under the generic reading,
$\forall x.\ F\,x\neq x$ is the disequality $F \neq \lambda x.x$: violated by $F := \lambda z.z$, satisfied by
$F:=\lambda z.c$. Under the extensional reading it says that $F\,t \neq t$ for \emph{every} $t$ of the
universe, so $F := \lambda z.c$ violates it ($t = c$); since it cannot be closed by an abstraction, it is printed
under a \scm{∀}.

\paragraph{Projection.} micro drops a disequality over a variable that is out of the answer, since such a
variable can always be chosen to keep the sides apart. That fails for a variable inside an argument of a flexible
term: $\exists Y.\ G \neq F\,Y$ is $F \neq \lambda z.G$, since $F$ may ignore its argument. Such an entry is kept,
and its variable prints as a working variable, read existentially:
\begin{scheme}
(°->pretty (fresh° (q) (fresh° (F G Y) (=° q (list F G)) (≠° G (@ F Y)))))
\end{scheme}
\begin{answers}
(((α β) : (≠ β (@ α _3))))
\end{answers}

\subsection{Absence: the \texorpdfstring{\texttt{parts}}{parts} hook}

\scm{absent°} looks for its datum in the \emph{parts} of a λ-term (\scm{μkanren-λ-parts}), never in binder names
or indices. A flexible term has no parts yet: it is \scm{'stuck} and watched whole, like a variable, and decided
once its head is known. Below a binder, the parts that mention the binder stay closed ($\lambda x.(a\,.\,b)$ is
looked into as $\lambda x.a$ and $\lambda x.b$), so a watched term never has a loose index.
\begin{scheme}
(°->list #f (fresh° (q F) (=° q F) (absent° 'c (@ F 'c))))
(°->list/datum (fresh° (q F) (=° q F) (absent° 'c (@ F 'c)) (=° F (ƛ (z) 'b))))
(°->list #f (fresh° (q) (absent° 'x (ƛ (x) `(g ,x))) (=° q 'ok)))
(°->list #f (fresh° (q F) (=° q F) (absent° 'c (ƛ (x) `(,(@ F x) c)))))
(°->list #f (fresh° (q F) (=° q F) (∀° (x) (absent° x (@ F x)))))
(°->list #f (fresh° (q) (∀° (x) (absent° x q))))
\end{scheme}
\begin{answers}
((λ (α) (assert (absent? 'c (@ α 'c))) α))
((ƛ (z) b))
((λ () 'ok))
()
((λ (α) (∀ (x) (assert (absent? x (@ α x)))) α))
((λ (α) α))
\end{answers}
In order: $c \notin F\,c$ waits; $F := \lambda z.b$ drops the argument, so it holds; the binder name \texttt{x}
is not the symbol \scm{'x}; $c$ is found at once outside the flexible part; an absento of an eigenvariable on a
term that may still hold it keeps $x$ and prints under \scm{∀}; and $\forall x.\ x \notin q$ is entailed by the
levels ($q$ has level $0<1$, so no instance of $q$ can hold $x$) and dropped. The datum may itself be a λ-term,
or reduce to an atom (\scm{(@ G 'b)} is \scm{a} once $G := \lambda z.a$): it is then unified, modulo βη, as micro
unifies a compound datum.

\subsection{Types: the \texorpdfstring{\texttt{tag}}{tag} hook}

A type on a flexible term waits in E as \scm{(tag tag u)}; the \texttt{post} hook does the same for a variable of
A just bound to one. A rigid λ-term or an eigenvariable is decided by the type's predicate: abstractions and
eigenvariables are neither symbols nor numbers, while a type of the user's made with \scm{define-μkanren-tag}
may hold of them. micro's types are disjoint, so two different types on the same flexible term fail at once:
\begin{scheme}
(°->list #f (fresh° (q F) (=° q (@ F 'a)) (symbol° q)))
(°->list/datum (fresh° (q F) (=° q (@ F 'a)) (symbol° q) (=° F (ƛ (z) z))))
(°->list #f (fresh° (q F) (=° q (@ F 'a)) (symbol° q) (number° q)))
(°->list #f (fresh° (q) (∀° (x) (symbol° x))))
(°->list #f (fresh° (q F) (=° q F) (∀° (x) (symbol° (@ F x)))))
\end{scheme}
\begin{answers}
((λ (α) (assert (every (μ v (symbol? v)) (list (@ α 'a)))) (@ α 'a)))
(a)
()
()
((λ (α) (∀ (x) (assert (every (μ v (symbol? v)) (list (@ α x))))) α))
\end{answers}

\begin{pitfall}[Readings are per constraint, and only for $\neq$]
The extensional reading applies to \scm{≠°} only; \scm{absent°} and the types always read an eigenvariable as
\rulename{generic} does. A disequality keeps the reading in force when it was posted; a first-order disequality that
migrates from D to N later takes the reading in force when it migrates, since D records none. The two can differ
only if the program changes \scm{μkanren-λ-reading} between the two goals.
\end{pitfall}

% ------------------------------------------------------------------------------------------------
\section{Reification}\label{sec:module-reify}

micro reifies an answer as a λ-expression whose parameters name the answer's variables and whose body lists the
constraints and ends with the value (\cref{sec:micro-reify}). The module plugs into that through the hooks
\texttt{value}, \texttt{repr}, \texttt{terms}, \texttt{forms} and \texttt{environment} (\cref{ch:design}).

\paragraph{Values.} Deep values are β-normal and η-short (\scm{μkanren-λ-value}). They are used by answers, by
D, A and T, by \scm{project°} and by aggregation.

\paragraph{Binders.} An abstraction prints as \scm{(ƛ (x y) body)}, nested binders merged. A binder is named
after its hint, renamed apart from the binders around it, from every symbol of its body, and from the operators
of the rebuilding code (\scm{cons}, \scm{quote}, \scm{vector}, \scm{make-record-instance}, \scm{ƛ}, \scm{@},
\scm{eigen}, \scm{bvar}). So neither capture nor shadowing can mislead the reader (Ma and Friedman's
\enquote{true = false}, \cref{sec:lkanren-merge}), and a ground answer evaluates back to the same term:
\begin{scheme}
(°->list/datum (fresh° (q) (=° q (ƛ (a) (ƛ (a) a)))))
(μkanren-datum (μkanren-abstract 'cons (λ (c) `(a ,c))))
\end{scheme}
\begin{answers}
((ƛ (a a1) a1))
(ƛ (cons1) (a cons1))
\end{answers}

\paragraph{Applications} print as \scm{(@ h a ...)}, so a flexible application \scm{(@ α x)} cannot be mistaken
for the list \scm{(α x)}.

\paragraph{Residual forms.} The constraints of the λ layer print as Scheme forms that can be evaluated:
\begin{fullwidth}
\medskip\small\centering
\begin{tabular}{@{}p{0.36\linewidth}p{0.58\linewidth}@{}}
\toprule
constraint & form in the answer \\
\midrule
suspended equation (E) & \texttt{(assert (≐ u v))} \\
type on a flexible term (E) & \texttt{(assert (every (μ v (symbol? v)) (list u)))} \\
absento with a pending term (E, T) & \texttt{(assert (absent? d u))} \\
disequality whose trial binds (N) & \texttt{(begin (deny (equal? α t)))}, with \texttt{(∃ (\_3) ...)} inside \texttt{deny} when the trial creates variables \\
disequality whose trial only suspends (N) & \texttt{(begin (deny (≐ u v)))} \\
any of them still holding an eigenvariable & wrapped in \texttt{(∀ (x) form)} \\
\bottomrule
\end{tabular}\par\bigskip
\end{fullwidth}
\scm{≐} is exported: it decides αβη-equality of closed terms by unifying them in the empty state, so the forms
evaluate. Through the \texttt{terms} hook, the variables of residual equations are reified \emph{after} those of
the answer, so they are named \texttt{α}, \texttt{β}, \dots{} rather than \texttt{\_2}
(\cref{ex:module-forget}). micro's show module gives the αKanren-style view \scm{°->pretty}:
\begin{scheme}
(°->pretty (fresh° (q) (∀° (x) (=° (@ q `(g ,x)) `(h (g ,x))))))
(°->pretty (fresh° (q) (∀° (x) (=° (@ q x x) x))))
\end{scheme}
\begin{answers}
((α : (= (ƛ (x) (@ α (g x))) (ƛ (x) (h (g x))))))
((α : (= (ƛ (x) (@ α x x)) (ƛ (x) x))))
\end{answers}

\paragraph{No free eigenvariable.} Every constraint that holds an eigenvariable past its scope is closed, by
\scm{ƛ} or by \scm{∀}, so no eigenvariable ever appears free in an answer. The one value that could hold one is the
answer itself, when the variable that \scm{°->list} reifies is created inside a \scm{∀°}. λProlog reports no binding
for a variable quantified inside \lp{pi x\ }, and this answer cannot be told without $x$, so reifying it is an error:
\begin{scheme}
(error-message (λ () (°->list #f (∀° (x) (fresh° (q) (=° q x))))))
\end{scheme}
\begin{answers}
"the answer holds the eigenvariable of a ∀° it was created in: create the variable it reifies outside every ∀°"
\end{answers}

\paragraph{Live ground answers.} \scm{ƛ}, \scm{@} and \scm{≐} are bound in the module's environment, which the
\texttt{environment} hook hands to \scm{°->list/ground}. A ground answer is therefore a λ-term one can apply:
\begin{scheme}
(@ (car (°->list/ground (fresh° (q) (=° q (ƛ (x y) `(g ,x ,y)))))) 'b 'b)
\end{scheme}
\begin{answers}
(g b b)
\end{answers}
\scm{°->list/datum} answers plain data instead, comparable with \scm{equal?} against quoted book answers; that is what
the suite uses, and \scm{μkanren-repr->datum} turns micro's rebuilding code into such data.

\begin{pitfall}[Hints are not part of equality]
\scm{equal?} sees hints, the unifier does not. \scm{(ƛ (x) x)} and \scm{(ƛ (y) y)} unify and are \scm{μkanren-α=?},
but they are not \scm{equal?}: \scm{equal?} compares records slot by slot, hint included, and the hint is
kept in the record so that answers print the user's binder names. When comparing λ-terms in host code, use \scm{μkanren-α=?} or \scm{≐}.
A related trap is Scheme's own: \scm{(ƛ (quote) 'a)} shadows \scm{quote} while the body is built; terms built through
\scm{μkanren-abstract} with such a hint print and evaluate correctly.
\end{pitfall}

% ------------------------------------------------------------------------------------------------
\section{API reference}\label{sec:module-api}

\begin{longtable}{@{}>{\raggedright\arraybackslash}p{0.40\fullwidthlen}>{\raggedright\arraybackslash}p{0.56\fullwidthlen}@{}}
\toprule
name & meaning \\
\midrule
\endfirsthead
\toprule
name & meaning \\
\midrule
\endhead
\multicolumn{2}{@{}l}{\emph{Terms}}\\
\scm{(ƛ (x ...) body)} & λ-abstraction over Scheme binders, η-short \\
\scm{(@ f a ...)} & application; β at once on an abstraction, nested applications flattened \\
\scm{(≐ u v)} & αβη-equality of closed terms, for residual forms \\
\scm{(μkanren-abstract name f)} & the abstraction of the term \scm{(f p)}, \scm{p} a placeholder; binder hinted \scm{name} \\
\scm{(μkanren-apply f args)} & \scm{@} on a list of arguments \\
\scm{(μkanren-hnorm t s)} & head normal form of \scm{t} in state \scm{s} (\cref{def:module-hnf}) \\
\scm{(μkanren-α=? u v)} & α-equivalence of deep values, hints ignored \\
\scm{(μkanren-datum t)} & a term as data, printed without a state \\
\scm{μkanren-lam?}\newline \scm{μkanren-app?}\newline \scm{μkanren-eigen?}\newline \scm{μkanren-bvar?}\newline \scm{μkanren-λ-term?} & record predicates; the last is their union \\
\scm{μkanren-lam-hint}\newline \scm{μkanren-eigen-name}\newline \scm{μkanren-eigen-level} & accessors \\
\multicolumn{2}{@{}l}{\emph{Goals}}\\
\scm{(∀° (x ...) g ...)} & \rulename{generic}; \scm{(∀ª name f)} its procedural form, \scm{f} taking the eigenvariable \\
\scm{(⊃° (D ...) g ...)} & \rulename{augment}; \scm{(⊃ª clauses g)} its procedural form \\
\scm{(define-dynamic-relation}\newline \scm{  (r a ...) g ...)} & a relation that hypotheses can extend; with no \scm{g}, no static clauses \\
\scm{(μkanren-dynamic f)}\newline \scm{(μkanren-dynamic? f)} & tag a procedure as dynamic, test the tag \\
\scm{μkanren-hyps°}\newline \scm{μkanren-clause}\newline \scm{μkanren-clause/vars}\newline \scm{μkanren-clause*} & what the macros expand to \\
\multicolumn{2}{@{}l}{\emph{Clauses $D$}}\\
\scm{(r t ...)}\newline \scm{((r t ...) ← g ...)}\newline \scm{(← (r t ...) g ...)}\newline \scm{(∀ (v ...) D)} & fact, rule (infix and prefix), own variables \\
\multicolumn{2}{@{}l}{\emph{Answers}}\\
\scm{(°->list/datum g)} & the answers of \scm{g} as plain data, constraints dropped \\
\scm{(μkanren-repr->datum r)} & micro's rebuilding code as data \\
\multicolumn{2}{@{}l}{\emph{Engine}}\\
\scm{μkanren-λ-reading} & parameter, \scm{'generic} (default) or \scm{'extensional} \\
\scm{(μkanren-λ-depth s)} & the current universe level \\
\scm{(μkanren-λ-level α s)} & the level of variable \scm{α} \\
\scm{(μkanren-λ-delayed s)} & the suspended constraints E \\
\scm{(μkanren-λ-disequalities s)} & the λ-disequalities N \\
\scm{μkanren-λ-extension} & the registered extension record \\
\multicolumn{2}{@{}l}{\emph{In \texttt{(aux kanren micro show)}}}\\
\scm{(°->pretty g)}\newline \scm{μkanren-answer->pretty} & αKanren-style answers, \scm{(value : constraint ...)} \\
\bottomrule
\end{longtable}

% ------------------------------------------------------------------------------------------------
\section{Summary}\label{sec:module-summary}

\begin{itemize}\emergencystretch=2em
  \item \textbf{Terms.} Four records, locally nameless: \texttt{μkanren-\allowbreak lam} (interned hint, body),
    \texttt{μkanren-\allowbreak app}, \texttt{μkanren-\allowbreak eigen} (name, level, uninterned id), \texttt{μkanren-\allowbreak bvar}. \scm{ƛ} uses
    Scheme binders and builds η-short; \scm{@} β-reduces at once. Terms built alike are \scm{equal?}; α-equality is
    \scm{μkanren-α=?} and \scm{≐}. Lists remain data: \scm{(@ 'g 'a)} is not \scm{'(g a)}.
  \item \textbf{State.} One immutable record in micro's slot \texttt{X}: depth, top, the levels $L$/$L_n$ (an
    sbral, lazily padded), Γ (one entry per relation), E (suspended constraints), N (λ-disequalities). \scm{∀°}
    and \scm{⊃°} restore depth and Γ on every answer; branches never share it.
  \item \textbf{Unification.} micro does the first-order part; the hook applies \rulename{ξ}, \rulename{η},
    \rulename{Rigid}, \rulename{FF-same}, \rulename{FR}/FF-diff by variable elimination, and suspends
    non-patterns: the rules of \cref{fig:rules-pattern} with the prefix erased into levels. Patterns are read
    off levels. Elim prunes and lowers only strict occurrences, fails an
    occurrence of $F$ only by Nipkow's size argument, and suspends whatever a later binding could decide.
  \item \textbf{Check.} Every binding is scope-checked at its variable's level (NJK's binding rule with pruning),
    or taken back into E when a non-pattern decides it. Fast paths skip it before any \scm{∀°}, on plain terms and
    for top-level variables. Wake-up retries E in rounds while a round binds, and terminates when hnf does.
  \item \textbf{∀° and ⊃°.} \rulename{generic} runs its body one universe deeper with a new eigenvariable and discharges
    the residuals on exit (raise, close by ξ, retry, repeat). \rulename{augment} pushes closures on Γ. Dynamic relations
    backchain on hypotheses first, most recent first, then on their static body; a clause's free variables are
    shared, its \scm{∀} variables renamed per use.
  \item \textbf{Stores.} \scm{≠°} on λ-terms goes to N with a trial whose new variables are existential;
    D never holds a λ-term; two readings of $\forall$, generic (default) and extensional. \scm{absent°} looks into
    parts and watches flexible terms; types wait on flexible terms and are decided by their predicate.
  \item \textbf{Answers.} β-normal, η-short, binders renamed apart; residuals as \scm{(assert (≐ u v))},
    \scm{(deny ...)}, \scm{(absent? ...)}, closed with \scm{∀} when they keep an eigenvariable; never a free
    eigenvariable; ground answers are live terms.
\end{itemize}
